\chapter{Periodic Structure of the I Ching System}
\label{chap:yijing-periodic}
\markboth{Periodic Structure of the I Ching System}{Periodic Structure of the I Ching System}

\begin{remark}[Periodicity decision for the target of this chapter (citation)]\label{rem:yijing_ps_hook}
When this chapter decides whether a research target is periodic, determines its period type, or seeks its minimal positive period, reference is always made to the `Periodic Systems' chapter:
for the periodicity criterion see \cref{thm:ps_periodicity_decision},
for the discrimination of period types see \cref{def:ps_period_type}, \cref{prop:ps_type_discrimination},
for the search for the minimal positive period see \cref{def:ps_minimal_period_algorithm}, \cref{thm:ps_minimal_period_correctness},
and for the construction of period series and epicycle groups from data see \cref{thm:ps_series_reconstruction}, \cref{def:ps_epicycle_group}.
\end{remark}

\begin{remark}[Positioning of this chapter: developed along the five-step interface of the `Periodic Systems' chapter]\label{rem:yijing_five_steps_positioning}
The periodic structure of the I Ching system studies the mathematical characterisation of periodicity and phase relations in the *I Ching* symbol system. This chapter does not establish a separate periodicity criterion; instead it connects the I Ching symbol system to the five-step interface prescribed by the `Periodic Systems' chapter:
\begin{enumerate}[label=(\arabic*)]
    \item \textbf{Decision}: reduce the symbol systems of Taiji, the Two Modes, the Four Symbols, the Eight Trigrams, the Sixty-Four Hexagrams, the Heavenly Stems, the Earthly Branches and the sexagenary cycle to discrete dynamical systems, cite \cref{thm:ps_periodicity_decision} to decide whether they are periodic, and determine the period type;
    \item \textbf{Construction}: from data such as hexagram images, stems and branches, solar terms and Na-Yin, construct the amplitude series, the frequency series and the phase series, and assemble them into the I Ching epicycle group;
    \item \textbf{Adjustment}: dynamically adjust the I Ching epicycle group by a unified ordering rule, the tensor structure and the decrease of the residual energy;
    \item \textbf{Convolution}: introduce amplitude convolution, phase convolution and frequency convolution on the I Ching epicycle group, and connect them with periodic convolution, the generation of the sexagenary cycle and the He Tu--Lo Shu structure;
    \item \textbf{Integration}: connect to the tool chains of the three chapters `Discrete Sampling and Intercalation of Periodic Signals', `Harmonic Analysis' and `Time Series Analysis', forming an I Ching periodicity diagnostic procedure.
\end{enumerate}
The mathematical tools used in this chapter include: cyclic groups, irreducible representations, tensor products, the discrete Fourier transform, finite-order actions, periodic convolution, phase operators, nested epicycle generation, He Tu complex functions and Lo Shu matrices.
\end{remark}

\section{Connection to the five-step interface of the `Periodic Systems' chapter}
\label{sec:yijing_five_steps}

\subsection{Step one: periodicity decision for an I Ching target}

\begin{definition}[I Ching target and decision input]\label{def:yijing_target}
An \textbf{I Ching target} is one of the following objects:
\begin{enumerate}[label=(\roman*)]
    \item a hexagram hierarchy system, whose state space is $\mathbb F_2^n$ and whose evolution map is the binary bottom-up ordering advance;
    \item a stem--branch sequence, whose state space is $\mathbb Z_{10}\times \mathbb Z_{12}$ and whose evolution map advances the Heavenly Stems and the Earthly Branches in synchrony;
    \item a sexagenary cycle sequence, whose state space is the cyclic group $C_{60}$ and whose evolution map is $+1\pmod{60}$;
    \item a solar-term or Na-Yin sequence, whose state space is a finite cyclic group $C_T$ and whose evolution map is $+1\pmod T$.
\end{enumerate}
The \textbf{decision input} of an I Ching target is the pair
\[
\mathcal D_{\mathrm{Y}}=(\mathcal S,\Phi),
\]
where $\mathcal S$ is the state space and $\Phi:\mathcal S\to\mathcal S$ is the evolution map induced by ordering advance, synchronous advance or cyclic translation. If $\mathcal S$ is a finite set, the target is called a \textbf{finite I Ching target}; if $\mathcal S$ is an infinite set, the target is called an \textbf{infinite I Ching target}.
\end{definition}

\begin{definition}[I Ching decision map]\label{def:yijing_periodization}
Let an I Ching target be given by the decision input $\mathcal D_{\mathrm{Y}}=(\mathcal S,\Phi)$. If the states of the target are generated by iteration of $\Phi$, that is, if there exists $x_0\in\mathcal S$ such that
\[
x_t=\Phi^t(x_0),\qquad t\in\mathbb Z,
\]
then $\Phi$ is called the \textbf{I Ching decision map} of the target. In this case the I Ching target reduces to a discrete dynamical system $(\mathcal S,\Phi)$, and the criteria of the `Operations and Periodicity' chapter and of the `Periodic Systems' chapter may be applied directly.
\end{definition}

\begin{theorem}[Periodicity decision theorem for I Ching targets]\label{thm:yijing_periodicity_decision}
Let the decision map of an I Ching target be $\Phi$, and let $(\mathcal S,\Phi)$ be a discrete dynamical system. Then:
\begin{enumerate}[label=(\Alph*)]
    \item \textbf{Sufficient decision}: if there exists $m\in\mathbb N^+$ such that $\Phi^m=\mathrm{id}$, then the target is a uniformly periodic I Ching system; every orbit is closed, and every orbit length divides $m$.
    \item \textbf{Equivalence of the decision}: the target is a uniformly periodic I Ching system if and only if $\operatorname{ord}(\Phi)<\infty$.
    \item \textbf{Pointwise criterion}: the orbit of the target at the state $x_0$ is closed if and only if $x_0$ is a periodic point; in that case the orbit length is the period at that point.
    \item \textbf{Non-omissibility}: the conjunction of a conserved quantity and a symmetry group cannot replace $\operatorname{ord}(\Phi)<\infty$; an almost periodic I Ching target does not belong to the periodic I Ching systems.
\end{enumerate}
\end{theorem}

\begin{proof}
The \cref{thm:yijing_periodicity_decision} of this chapter is the domain instantiation of \cref{thm:ps_periodicity_decision} of the `Periodic Systems' chapter:
one need only replace the decision map $f$ and the evolution group by the $\Phi$ of this chapter, and the four items (A) sufficient decision, (B) equivalence of the decision, (C) pointwise criterion
and (D) non-omissibility hold one by one. By the convention of \cref{rem:ps_applications_hook}, they are not expanded again here.
\end{proof}

\begin{corollary}[Period types of the I Ching]\label{cor:yijing_period_type}
After the decision by \cref{thm:yijing_periodicity_decision}, an I Ching target falls into one of the following types:
\begin{enumerate}[label=(\roman*)]
    \item \textbf{Trivially periodic type}: the least period is $m_0=1$ and the state is a fixed point;
    \item \textbf{Strictly periodic type}: $1<m_0<\infty$ and the state is strictly periodic;
    \item \textbf{Uniformly periodic type}: $\operatorname{ord}(\Phi)=m<\infty$ and all orbit lengths divide $m$;
    \item \textbf{Quasi-periodic type}: $\operatorname{ord}(\Phi)=\infty$; the orbit does not close but is dense on the compact carrier structure;
    \item \textbf{Eventually periodic type}: there exists $N$ such that the target is periodic for $t\ge N$;
    \item \textbf{Discretely periodic type}: the carrier structure is $\mathbb Z$ and the normalised period ratio $\rho=T_x/\Delta x\in\mathbb Q$.
    \item \textbf{Asymptotically periodic (damped) type}: the orbit is not closed, but tends asymptotically to a closed orbit or to a fixed point, the residual decaying as $e^{-\kappa t}$ (see \cref{def:ps_period_type}(vii)).
\end{enumerate}
Among them (i)--(iii) and (vi) belong to the periodic I Ching systems; (iv) does not belong to the periodic I Ching systems, but can be handled by bounded approximation tools; (v) is a periodic I Ching system only in its stable stage; (vii), like (iv), is a type in the limiting sense; it is not a periodic I Ching system and can be handled with bounded-approximation tools only in its stable stage.
\end{corollary}

\begin{proof}
Each item follows by substituting an I Ching target into \cref{thm:yijing_periodicity_decision} and \cref{def:ps_period_type} of the `Periodic Systems' chapter.
\end{proof}

\begin{lemma}[Finite-order decision for a hexagram hierarchy]\label{lem:yijing_bagua_finite_order}
Let the state space of a hexagram hierarchy be $\mathbb F_2^n$ and let the evolution map be binary addition of one,
\[
\Phi(b_0,\dots,b_{n-1})=(b_0,\dots,b_{n-1})+1\pmod{2^n}.
\]
Then $\operatorname{ord}(\Phi)=2^n$, and the hexagram hierarchy is a uniformly periodic system.
\end{lemma}

\begin{proof}
Under addition the state space $\mathbb F_2^n$ is isomorphic to the cyclic group $C_{2^n}$. The map $\Phi$ is $+1$, whose order is the group order $2^n$. Hence $\Phi^{2^n}=\mathrm{id}$, and for every $k<2^n$ we have $\Phi^k\neq\mathrm{id}$. Therefore $\operatorname{ord}(\Phi)=2^n$. By \cref{thm:yijing_periodicity_decision}, the system is a uniformly periodic system.
\end{proof}

\begin{lemma}[Periodicity decision for the Heavenly Stems and Earthly Branches]\label{lem:yijing_ganzhi_finite_order}
The state space of the Heavenly Stems system is $C_{10}$ and the state space of the Earthly Branches system is $C_{12}$. Then
\[
\operatorname{ord}(\Phi_{\text{Heavenly Stems}})=10,\qquad
\operatorname{ord}(\Phi_{\text{Earthly Branches}})=12.
\]
Both are uniformly periodic systems.
\end{lemma}

\begin{proof}
The stem sequence is the cycle of the ten Heavenly Stems and the branch sequence is the cycle of the twelve Earthly Branches. By the definition of a cyclic group, $+1\pmod{10}$ has order $10$ and $+1\pmod{12}$ has order $12$. The conclusion follows.
\end{proof}

\begin{lemma}[Periodicity decision for the sexagenary cycle]\label{lem:yijing_ganzhi60_finite_order}
The state space of the sexagenary system is $C_{60}$ and the evolution map is $+1\pmod{60}$. Then
\[
\operatorname{ord}(\Phi_{\text{sexagenary cycle}})=60.
\]
The sexagenary cycle is a uniformly periodic system.
\end{lemma}

\begin{proof}
The sexagenary cycle is generated by the synchronous parity of the Heavenly Stems and the Earthly Branches. By \cref{ch20_ganzhi_parity_filter}, there are exactly $60$ effective combinations, and they are isomorphic to $C_{60}$. Hence $+1\pmod{60}$ has order $60$.
\end{proof}

\begin{definition}[Algorithm for the least positive period of an I Ching target]\label{def:yijing_minimal_period_algorithm}
This definition has the same content as Chapter~\ref{chap:periodic_system}, the section `Periodicity decision: is the target periodic?' (\cref{sec:ps_periodicity_decision}), \cref{def:ps_minimal_period_algorithm}; to avoid repeating the statement, only the reference anchor is retained here, and the complete statement is given there. The only difference here is terminology (I Ching target / I Ching sample); the algorithm steps and the conclusion are exactly the same.
\end{definition}

\begin{theorem}[Correctness of the algorithm for the least positive period of an I Ching target]\label{thm:yijing_minimal_period_correctness}
This theorem has the same content as Chapter~\ref{chap:periodic_system}, the section `Periodicity decision: is the target periodic?' (\cref{sec:ps_periodicity_decision}), \cref{thm:ps_minimal_period_correctness}; to avoid repeating the statement, only the reference anchor is retained here, and the complete statement and proof are given there. The only difference here is terminology (I Ching target / I Ching sample); the algorithm steps and the conclusion are exactly the same.
\end{theorem}

\begin{example}[Periodicity decision from Taiji to the Sixty-Four Hexagrams]\label{ex:yijing_taiji_to_64}
\begin{enumerate}[label=(\roman*)]
    \item Taiji: $T=1$, trivially periodic type;
    \item the Two Modes: $T=2$, strictly periodic type;
    \item the Four Symbols: $T=4$, uniformly periodic type;
    \item the Eight Trigrams: $T=8$, uniformly periodic type;
    \item the Sixty-Four Hexagrams: $T=64$, uniformly periodic type;
    \item the Heavenly Stems: $T=10$, uniformly periodic type;
    \item the Earthly Branches: $T=12$, uniformly periodic type;
    \item the sexagenary cycle: $T=60$, uniformly periodic type.
\end{enumerate}
\end{example}

\begin{remark}[An invariant does not guarantee periodicity]\label{rem:yijing_invariant_not_period}
Certain conserved quantities in the I Ching symbol system (such as parity synchronisation and the assignment of the Five Phases) are auxiliary conditions for periodic structure, but they do not constitute a sufficient condition. Periodic structure must be established by \cref{thm:yijing_periodicity_decision} or by the finite-order decision lemmas above.
\end{remark}

\subsection{Step two: construction of the I Ching period series and epicycle group}

\begin{definition}[I Ching amplitude series, frequency series and phase series]\label{def:yijing_series}
Let the state or output sequence of an I Ching target be $\{z_t\}$. If it can be written as
\[
z_t=\sum_k A_k e^{i(\omega_k t+\phi_k)}+\varepsilon_t,
\]
then
\[
\mathcal A_z(t)=\sum_k A_k e^{ikt},\qquad
\mathcal F_z(s)=\sum_k \omega_k s^k,\qquad
\mathcal P_z(t)=\sum_k e^{i(kt+\phi_k)}
\]
are called respectively the \textbf{amplitude series}, the \textbf{frequency series} and the \textbf{phase series} of the I Ching target. The three agree with \cref{def:amplitude_series}, \cref{def:frequency_series} and \cref{def:phase_series} of the `Series and Convolution' chapter.
\end{definition}

\begin{definition}[I Ching epicycle group]\label{def:yijing_epicycle_group}
Let the period series of an I Ching target be given by the amplitude series $\{A_k\}$, the frequency series $\{\omega_k\}$ and the phase series $\{\phi_k\}$. For each $k$ define the epicycle
\[
E_k:\theta\mapsto A_k e^{i(\theta+\phi_k)},\qquad \theta=\omega_k t,
\]
and call the set
\[
\mathfrak E_{\mathrm{Y}}=\{(A_k,\omega_k,\phi_k)\}_{k\in\mathbb Z}
\]
the \textbf{epicycle group} of the I Ching target. Its superposition is the I Ching period series
\[
z_t=\sum_k E_k(\omega_k t)=\sum_k A_k e^{i(\omega_k t+\phi_k)}.
\]
\end{definition}

\begin{theorem}[Construction of the period series from I Ching data]\label{thm:yijing_series_reconstruction}
Let the output of an I Ching target be a finite periodic sequence of period $N$. Then:
\begin{enumerate}[label=(\roman*)]
    \item the output can be expanded as a discrete Fourier series
    \[
    z_t=\sum_{k=0}^{N-1}c_k e^{i2\pi kt/N},\qquad
    c_k=\frac1N\sum_{t=0}^{N-1}z_t e^{-i2\pi kt/N};
    \]
    \item the amplitude series is given by $|c_k|$, the phase series by $\arg c_k$, and the frequency series by $2\pi k/N$;
    \item if the output is a hexagram sequence or a stem--branch sequence, then the harmonics at the various levels correspond to different period scales.
\end{enumerate}
\end{theorem}

\begin{proof}
(i) The discrete Fourier series of a finite periodic sequence is a standard result: the basis functions $\{e^{i2\pi kt/N}\}_{k=0}^{N-1}$ form an orthogonal basis of $\mathbb C^N$.
(ii) This is read off directly from $c_k=A_k e^{i\phi_k}$.
(iii) For a hexagram sequence $N=64$, the fundamental frequency corresponds to the great period of the Sixty-Four Hexagrams, and the lower harmonics correspond to substructures such as Yin-Yang, the Four Symbols and the Eight Trigrams; for a stem--branch sequence $N=60$, the fundamental frequency corresponds to the sexagenary period, and the harmonics correspond to the Heavenly Stems, the Earthly Branches and parity synchronisation.
\end{proof}

\begin{proposition}[Correspondence of the three parameters: radius, phase and angular speed]\label{prop:yijing_epicycle_correspondence}
The three parameters of the I Ching epicycle group $\mathfrak E_{\mathrm{Y}}=\{(A_k,\omega_k,\phi_k)\}_{k}$ correspond one-to-one with three geometric--mechanical quantities:
\begin{enumerate}[label=(\roman*)]
    \item \textbf{Radius corresponds to amplitude}: the radius $|E_k|=A_k$ of $E_k$ is the amplitude of that component;
    \item \textbf{Phase corresponds to rotational delay}: the initial phase $\phi_k$ corresponds to the time translation $\tau_k=\phi_k/\omega_k$, that is, the epicycle \textbf{leads} the reference by $\tau_k$ ($\phi_k>0$ means that the waveform is \textbf{advanced} by $\tau_k$ relative to the reference, equivalently the waveform coincides with the reference once the reference has been delayed by $\tau_k$);
    \item \textbf{Frequency corresponds to angular speed}: $\omega_k$ is the angular velocity of that epicycle about the origin.
\end{enumerate}
\end{proposition}

\begin{proof}
The \cref{prop:yijing_epicycle_correspondence} of this chapter corresponds exactly to \cref{prop:ps_epicycle_correspondence} of the `Periodic Systems' chapter;
the three correspondences---radius to amplitude, phase to rotational delay and frequency to angular speed---are obtained one by one by translating them to the epicycle group of this chapter.
The proof is to be found there: (i) from $|E_k|=A_k$ and the energy decomposition, (ii) from the phase-shift theorem, (iii) from the definition of the frequency series.
\end{proof}

\begin{example}[Period series of the Heavenly Stems]\label{ex:yijing_tiangan_series}
The ten Heavenly Stems may be regarded as a periodic sequence on $C_{10}$; its discrete Fourier decomposition yields ten frequency components, the fundamental frequency corresponding to the cycle of the ten Heavenly Stems, the second harmonic to the alternation of Yin-Yang and the fifth harmonic to the cycle of the Five Phases.
\end{example}

\begin{remark}[A period series is a complex function]\label{rem:yijing_series_complex}
Let $z=e^{i\omega_0t}$. Then the I Ching period series
\[
z_t=\sum_k A_k e^{i(\omega_k t+\phi_k)}
\]
can be written as the Laurent series $\sum_k c_k z^k$, where $c_k=A_k e^{i\phi_k}$, provided that \textbf{every component frequency is commensurable with $\omega_0$} (that is, $\omega_k=k\omega_0$ with $k\in\mathbb Z$). Hence \textbf{when all components are commensurable with $\omega_0$, the period series is a complex-valued function on the unit circle (otherwise it is a function on the torus)}.
\end{remark}

\begin{definition}[Unified ordering rule]\label{def:unified_ordering}
\begin{enumerate}[label=(\roman*)]
    \item $k=0$ ($\theta=0$) corresponds to the \textbf{beginning} of the period. When $T=2^n$ this is the all-yin line configuration (Zhen, the beginning of wood); for other periods it is \textbf{yang wood}.
    \item The intermediate phases are arranged in increasing order of $\theta_k = 2\pi k/T$.
    \item $k=T-1$ corresponds to the \textbf{end} of the period. When $T=2^n$ this is the all-yang line configuration (Kan, the end of water); for other periods it is \textbf{yin water}.
\end{enumerate}
For $T=2^n$, the lines are read in binary from the bottom upwards: the first line is $b_0$, a yin line is $=0$ and a yang line is $=1$.
\end{definition}

\subsection{Step three: dynamic adjustment of the I Ching epicycle group}

\begin{definition}[Dynamic adjustment of the I Ching epicycle group]\label{def:yijing_adaptive_adjustment}
Let the current I Ching epicycle group be $\mathfrak E^{(j)}$ and let the partial sum be
\[
z^{(j)}(t)=\sum_{k\in\mathfrak E^{(j)}}A_k e^{i(\omega_k t+\phi_k)},
\]
with residual $r^{(j)}=z-z^{(j)}$. The \textbf{dynamic adjustment} consists of two rules:
\begin{enumerate}[label=(\roman*)]
    \item \textbf{Cardinality adjustment}: if the residual energy $\|r^{(j)}\|^2>\delta$, add one epicycle according to the peak of the residual periodogram;
    \item \textbf{Parameter adjustment}: update the parameters of the existing epicycles by regularised least squares or by the unified ordering rule.
\end{enumerate}
In the I Ching system, cardinality adjustment corresponds to adding a hexagram level, a stem--branch level or a Na-Yin level; parameter adjustment corresponds to adjusting the tensor index, the Five-Phase assignment or the phase origin.
\end{definition}

\begin{proposition}[Monotonicity of the cardinality adjustment]\label{prop:yijing_adaptive_growth}
After adding the dominant frequency component of the residual as in \cref{def:yijing_adaptive_adjustment}(i), the residual energy is monotonically non-increasing:
\[
\|r^{(j+1)}\|^2\le\|r^{(j)}\|^2,
\]
with equality if and only if the residual has no energy at the selected frequency.
\end{proposition}

\begin{proof}
The \cref{prop:yijing_adaptive_growth} of this chapter agrees with the dynamic adjustment mechanism of \cref{def:ps_adaptive_adjustment} of the `Periodic Systems' chapter:
adding an epicycle removes the component of $r^{(j)}$ along $\hat\omega$ from the residual, and the orthogonal decomposition
$\|r^{(j)}\|^2=\|r^{(j+1)}\|^2+\|\text{the removed component}\|^2$ gives the inequality,
with equality if and only if the norm of the removed component is zero. The complete argument is to be found there.
\end{proof}

\begin{theorem}[Convergence of the dynamic adjustment]\label{thm:yijing_adaptive_convergence}
Let the I Ching target be a \textbf{real-valued finite periodic sequence} of length $T_0$ ($T_0\in\mathbb N$) and let $\delta=0$. Then after at most $\lfloor T_0/2\rfloor+1$ steps of dynamic adjustment the residual is zero, and the resulting I Ching epicycle group is the exact Fourier data of the target; for a \textbf{complex-valued} sequence the corresponding number of steps is at most $T_0$.
\end{theorem}

\begin{proof}
The discrete Fourier series of a \textbf{real-valued} finite periodic sequence of length $T_0$ ($T_0\in\mathbb N$) is complete on $\lfloor T_0/2\rfloor+1$ harmonics (including the constant term); a \textbf{complex-valued} sequence requires $T_0$ components. At each step one dominant frequency component of the residual is extracted as in \cref{prop:yijing_adaptive_growth} and the residual energy decreases strictly, so after finitely many steps the residual is zero.
\end{proof}

\begin{definition}[Structural-metric adjustment]\label{def:yijing_structure_metric_adjustment}
Let the structural metric of an I Ching target be
\[
\mathcal I=(\mathcal I_{\text{hexagram images}},\mathcal I_{\text{stems and branches}},\mathcal I_{\text{Five Phases}})
=(\text{product of bit widths},\ \text{degree of parity synchronisation},\ \text{closure of the Five-Phase assignment}).
\]
\textbf{Structural-metric adjustment} is: at each step choose a neighbouring parameter $\theta$ in the parameter space $\Theta$ for which $\mathcal I(\theta)$ is \textbf{strictly decreasing} in the lexicographic order (minimisation); if no such neighbouring parameter exists, the process stops---then $\theta$ is a \textbf{locally optimal} parameter.
\end{definition}

\begin{proposition}[Finite convergence of the structural-metric adjustment]\label{prop:yijing_structure_metric_convergence}
Let the parameter space $\Theta$ be finite and let each adjustment make $\mathcal I$ \textbf{strictly decreasing} in the lexicographic order. Then the adjustment terminates after finitely many steps at a locally optimal parameter.
\end{proposition}

\begin{proof}
$\Theta$ is finite, so $\mathcal I$ takes finitely many values. The lexicographic order is well-founded on a finite set of values, so a strictly lexicographically decreasing sequence cannot continue indefinitely and the process terminates after finitely many steps. At termination there is no neighbouring parameter making $\mathcal I$ strictly lexicographically decreasing, that is, the parameter is locally optimal. Monotonicity must be imposed on the \textbf{whole lexicographic order}: if only `some component strictly increases or decreases' is required, cycles such as $(2,1)\to(1,2)\to(2,1)$ cannot be excluded and finite termination fails.
\end{proof}

\begin{remark}[Boundaries of the dynamic adjustment]\label{rem:yijing_adaptive_boundary}
\cref{thm:yijing_adaptive_convergence} relies on a finite periodic sequence. For a continuous or almost periodic I Ching target there is no exact representation by a finite epicycle group, and only bounded approximation is possible. Structural-metric adjustment does not guarantee a global optimum.
\end{remark}

\subsection{Step four: I Ching convolution}

\begin{definition}[Phase-operator periodic convolution]\label{def:PPO_periodic_convolution}
Let two symbol systems $A$ and $B$ have periods $T_A$ and $T_B$ and phase-operator sets $\mathcal{P}_{T_A}$ and $\mathcal{P}_{T_B}$ respectively. The \textbf{phase-operator periodic convolution} of $A$ and $B$ is defined by
\[
(\mathcal{P}_{T_A} \circledast \mathcal{P}_{T_B})(n, m) = \mathcal{O}_n^{(T_A)} \cdot \mathcal{O}_m^{(T_B)} = e^{i2\pi(n/T_A + m/T_B)}.
\]
The periodic convolution retains the interaction information between all pairs $(n,m)$---it is the complete set of all phase-relation operators between the two periodic systems.
\end{definition}

\begin{proposition}[Consistency of periodic convolution with nested epicycle generation]\label{prop:periodic_conv_epicycle_consistency}
Nested epicycle generation is the realisation of periodic convolution in a nested epicyclic structure. The composite phase operator $\mathcal{O}_{a,b} = \mathcal{O}_a^{(T_{\text{inner}})} \cdot \mathcal{O}_b^{(T_{\text{outer}})}$ is precisely the $(a,b)$ element of the periodic convolution matrix.
\end{proposition}

\begin{proof}
Take $A$ to be the inner-circle system ($T_A = T_{\text{inner}}$) and $B$ the outer-circle system ($T_B = T_{\text{outer}}$). By definition \ref{def:PPO_periodic_convolution}, the periodic convolution takes at $(a,b)$ the value
\[
(\mathcal{P}_{T_{\text{inner}}} \circledast \mathcal{P}_{T_{\text{outer}}})(a,b) = \mathcal{O}_a^{(T_{\text{inner}})} \cdot \mathcal{O}_b^{(T_{\text{outer}})} = e^{i2\pi(a/T_{\text{inner}} + b/T_{\text{outer}})}.
\]
By theorem \ref{thm:epicycle_nesting_generation}, the composite phase operator $\mathcal{O}_{a,b}$ of nested epicycle generation is exactly the same expression. The two are equal element by element.
\end{proof}

\begin{theorem}[Periodic-convolution generation of the sexagenary cycle]\label{thm:ganzhi_periodic_convolution}
The sexagenary cycle is generated by the periodic convolution of the Heavenly Stems ($T=10$) and the Earthly Branches ($T=12$). The effective-combination condition $n \equiv m \pmod{2}$ selects $60$ effective phases.
\end{theorem}

\begin{proof}
By definition \ref{def:PPO_periodic_convolution}, the periodic convolution of the Heavenly Stems ($T_A = 10$) and the Earthly Branches ($T_B = 12$) runs through all $10 \times 12 = 120$ pairs $(n,m)$, giving $120$ composite phase operators $\mathcal{O}_n^{(10)} \mathcal{O}_m^{(12)}$. Yin-Yang synchronisation requires the parities to agree, so the effective-combination condition is $n \equiv m \pmod 2$. By lemma \ref{ch20_ganzhi_parity_filter}, this condition selects $60$ of the $120$ pairs, and these correspond one-to-one with the sexagenary cycle; likewise, by the group structure proved in that lemma, these $60$ phases constitute the cyclic group $C_{60}$, so the period of the sexagenary cycle is $T = 60$.
\end{proof}

\begin{corollary}[Uniformity of the periodic convolution over all periodic systems]\label{cor:all_periods_periodic_conv}
All phase relations between any two periodic systems can be computed uniformly by means of the phase-operator periodic convolution.
\end{corollary}

\begin{proof}
This is a direct corollary of definition \ref{def:PPO_periodic_convolution}: by definition the periodic convolution $\mathcal{O}_n^{(T_A)} \mathcal{O}_m^{(T_B)}$ runs through all pairs $(n,m)$, so every phase-relation operator between the two periodic systems lies in the periodic convolution matrix, with nothing omitted; and every element of the matrix is in turn a product of two phase operators, which by the closure property of proposition \ref{prop:PPO_relation_structure} is still a phase operator, so there is nothing beyond it.
\end{proof}

\subsection{Step five: integration with discrete sampling, harmonic analysis and time series analysis}

\begin{definition}[I Ching periodicity diagnostic procedure]\label{def:yijing_periodic_diagnosis}
For a given I Ching symbol system, the diagnosis is carried out using the framework of periodic mathematics:
\begin{enumerate}[label=(\roman*)]
    \item \textbf{Periodicity decision}: verify the period type of the I Ching target (\cref{thm:yijing_periodicity_decision});
    \item \textbf{Series construction}: from hexagram, stem--branch and solar-term data construct the amplitude, frequency and phase series and the epicycle group (\cref{thm:yijing_series_reconstruction});
    \item \textbf{Tensor diagnosis}: compute the prime factorisation, the group structure, the tensor shape and the source of nested generation;
    \item \textbf{Convolution diagnosis}: compute the periodic convolution, the generation of the sexagenary cycle and the He Tu--Lo Shu structure;
    \item \textbf{Integrated diagnosis}: connect to the tool chains of discrete sampling, harmonic analysis and time series analysis.
\end{enumerate}
\end{definition}

\begin{theorem}[The I Ching periodicity diagnostic procedure and the interface with three chapters]\label{thm:yijing_diagnosis_integration}
The five steps of \cref{def:yijing_periodic_diagnosis} connect respectively to:
\begin{enumerate}[label=(\roman*)]
    \item the chapter `Discrete Sampling and Intercalation of Periodic Signals', \cref{def:uniform_sampling_parallel}, \cref{thm:nyquist_shannon_parallel}, \cref{thm:intercalation_leakage_suppression};
    \item the chapter `Harmonic Analysis', \cref{def:fourier_series}, \cref{thm:parseval_fourier}, \cref{prop:twelve_to_six};
    \item the chapter `Time Series Analysis', \cref{def:ts_period_detection}, \cref{prop:ts_period_significance}, \cref{def:ts_window_validation}.
\end{enumerate}
Hence the I Ching periodicity diagnosis is not an independent method, but an application of the unified tool chain of periodic mathematics to a traditional symbol system.
\end{theorem}

\begin{proof}
The \cref{thm:yijing_diagnosis_integration} of this chapter is the implementation of the five-step interface of the `Periodic Systems' chapter (\cref{tab:ps_five_steps}) in this domain:
the periodicity decision uses sampling and spectral estimation, the frequency decomposition and the phase kernel use the `Harmonic Analysis' chapter, and the periodogram peak, Fisher's test and
window-function validation use the `Time Series Analysis' chapter. Hence the diagnostic procedure is composed of the tool chains of three chapters, and the proof is to be found there.
\end{proof}

\begin{remark}[Operational order after integration]\label{rem:yijing_integration_order}
The recommended order for an actual diagnosis is: first decide whether the I Ching target is periodic (\cref{thm:yijing_periodicity_decision}); then construct the Fourier data and the epicycle group of the hexagram, stem--branch and solar-term data (\cref{thm:yijing_series_reconstruction}); then optimise the parameters by the unified ordering rule and the dynamic adjustment (\cref{thm:yijing_adaptive_convergence}, \cref{prop:yijing_structure_metric_convergence}); then carry out the analysis of the periodic convolution and of sexagenary generation (\cref{thm:ganzhi_periodic_convolution}); and finally output the diagnosis in accordance with \cref{def:yijing_periodic_diagnosis}.
\end{remark}

\section{Periodic phase operators}

\begin{definition}[Periodic phase operator]\label{def:PPO_yijing}
For any period $T \in \mathbb{N}^+$, the complex number $z_k^{(T)} = e^{i2\pi k/T}$ corresponding to its $k$-th phase is called the \textbf{periodic phase operator} and is written $\mathcal{O}_k^{(T)}$. The $T$ phase operators together form the set:
\[
\mathcal{P}_T = \{\mathcal{O}_k^{(T)} = e^{i2\pi k/T} \mid k = 0, 1, \dots, T-1\}.
\]

A periodic phase operator has a \textbf{dual identity}:
\begin{enumerate}[label=(\roman*)]
    \item \textbf{As a vector}: it is the direction vector on the unit circle of the $k$-th phase point in the period $T$---it encodes the position of that phase within the period.
    \item \textbf{As an operator}: it acts on any complex vector by complex multiplication, producing a phase rotation:
    \[
    \mathcal{O}_k^{(T)} \cdot \mathbf{v} = \mathbf{v} \cdot e^{i2\pi k/T}.
    \]
    Multiplication by $\mathcal{O}_k^{(T)}$ is equivalent to rotating the vector $\mathbf{v}$ through $2\pi k/T$ radians.
\end{enumerate}
\end{definition}

\begin{proposition}[Relational structure of the phase operators]\label{prop:PPO_relation_structure}
The $T$ phase operators of period $T$ form a complete relational hierarchy:
\begin{enumerate}[label=(\roman*)]
    \item \textbf{Yang operator (identity)}: $\mathcal{O}_0^{(T)} = 1$. Multiplication by $\mathcal{O}_0$ gives the object itself---the identity relation.
    \item \textbf{Yin operator (opposition)}: $\mathcal{O}_{T/2}^{(T)} = -1$ (when $T$ is even). Multiplication by $\mathcal{O}_{T/2}$ gives the opposite vector---the opposition relation.
    \item \textbf{General rotation operator}: $\mathcal{O}_k^{(T)}$ ($k \neq 0, T/2$). Partial rotation relations.
\end{enumerate}
All phase operators together form the cyclic group $C_T$ under complex multiplication: $\mathcal{O}_a \cdot \mathcal{O}_b = \mathcal{O}_{a+b \bmod T}$.
\end{proposition}

\begin{proof}
(i) Yang operator: $\mathcal{O}_0^{(T)} = e^{i\cdot 0} = 1$, so $\mathcal{O}_0^{(T)} z = z$ for every $z$. (ii) Yin operator: when $T$ is even, $T/2 \in \mathbb{N}$, so $\mathcal{O}_{T/2}^{(T)} = e^{i2\pi(T/2)/T} = e^{i\pi} = -1$. (iii) General rotation: for any $z = |z|e^{i\varphi}$, $\mathcal{O}_k^{(T)} z = |z| e^{i(\varphi + 2\pi k/T)}$; the modulus is unchanged and the argument increases by $2\pi k/T$.

Verification of the group structure: for any $a,b$, by the addition theorem for the exponential function,
\[
\mathcal{O}_a^{(T)} \cdot \mathcal{O}_b^{(T)} = e^{i2\pi a/T} e^{i2\pi b/T} = e^{i2\pi(a+b)/T} = \mathcal{O}_{a+b \bmod T}^{(T)}.
\]
Hence $\mathcal{P}_T$ is closed under complex multiplication; associativity is inherited from multiplication of complex numbers; the identity element is $\mathcal{O}_0^{(T)} = 1$; and the inverse of $\mathcal{O}_a^{(T)}$ is $\mathcal{O}_{T-a \bmod T}^{(T)}$. Therefore $(\mathcal{P}_T, \cdot)$ is a cyclic group, generated by $\mathcal{O}_1^{(T)}$, and $\mathcal{O}_1^{(T)}$ has order $T$, so $\mathcal{P}_T \cong C_T$.
\end{proof}

\begin{proposition}[Correspondence between phase operators and irreducible representations]\label{prop:PPO_group_yijing}
The set $\mathcal{P}_T$ forms the cyclic group $C_T$ under complex multiplication, and corresponds one-to-one with the irreducible representations of $C_T$:
\[
\mathcal{O}_k^{(T)} = \rho_k(g).
\]
\end{proposition}

\begin{proof}
By proposition \ref{prop:PPO_relation_structure}, $(\mathcal{P}_T, \cdot) \cong C_T$ with isomorphism $\chi: C_T \to \mathcal{P}_T,\; g^j \mapsto \mathcal{O}_j^{(T)}$. Each one-dimensional irreducible representation $\rho_k$ ($k = 0,\dots,T-1$) of the cyclic group $C_T = \langle g \mid g^T = 1 \rangle$ is uniquely determined by the image of the generator $g$: $\rho_k(g) = \omega^k$, where $\omega = e^{i2\pi/T}$. Thus $\rho_k(g^j) = \omega^{kj} = \mathcal{O}_{kj \bmod T}^{(T)}$, that is, the value of the irreducible representation $\rho_k$ at the generator is precisely the phase operator $\mathcal{O}_k^{(T)}$. Since $C_T$ has $T$ conjugacy classes, it has exactly $T$ one-dimensional irreducible representations, and both $k \mapsto \rho_k$ and $k \mapsto \mathcal{O}_k^{(T)}$ are bijections, so $\mathcal{P}_T$ corresponds one-to-one with the irreducible representations of $C_T$.
\end{proof}

\begin{definition}[Regular representation and DFT diagonalisation]\label{def:regular_representation_yijing}
This definition has the same content as Chapter~\ref{chap:harmonic}, the section `Roots of unity and cyclic group representations' (\cref{sec:unit_roots_cyclic_rep}), \cref{def:regular_representation_dft}; to avoid repeating the statement, only the reference anchor is retained here, and the complete statement is given there. Here the definition is merely restated in the language of the phase operators of the cyclic group $\mathbb{Z}_T$; the mathematical object is the same.
\end{definition}

\begin{remark}[The relational essence of the phase operators]
The core significance of phase operators lies in \textbf{defining relations}. The composition $\mathcal{O}_a \cdot \mathcal{O}_b$ of any two phases $\mathcal{O}_a$ and $\mathcal{O}_b$ gives a relation from $\mathcal{O}_b$ to $\mathcal{O}_{a+b}$. The multiplicative closure of the phase operators guarantees the completeness of the relational system.
\end{remark}

\section{Nested epicycle generation}

\begin{definition}[Nested epicycle generation rule]\label{def:epicycle_nesting_rule}
Let an epicycle of period $T_{\text{inner}}$ (the inner circle) be nested on an epicycle of period $T_{\text{outer}}$ (the outer circle). The \textbf{nested epicycle generation rule} is:
\[
\text{For every single step forward of the outer circle (a rotation of } 2\pi/T_{\text{outer}} \text{), the inner circle rotates through one full turn (a rotation of } 2\pi \text{).}
\]
Equivalently, the angular velocity of the inner circle is $T_{\text{inner}}$ times that of the outer circle: $\omega_{\text{inner}} = T_{\text{inner}} \cdot \omega_{\text{outer}}$. The period of the composite system is $T_{\text{composite}} = T_{\text{inner}} \times T_{\text{outer}}$.
\end{definition}

\begin{theorem}[Nested epicycle generation theorem]\label{thm:epicycle_nesting_generation}
Let the inner epicycle have $T_{\text{inner}}$ equally spaced phase points and the outer epicycle have $T_{\text{outer}}$ equally spaced phase points. The phases of the composite system are
\[
\theta_{a,b} = \frac{2\pi a}{T_{\text{inner}}} + \frac{2\pi b}{T_{\text{outer}}}, \quad a = 0, \dots, T_{\text{inner}}-1,\; b = 0, \dots, T_{\text{outer}}-1.
\]
The composite system has $T_{\text{inner}} \times T_{\text{outer}}$ independent phases, and its composite phase operator is $\mathcal{O}_{a,b} = \mathcal{O}_a^{(T_{\text{inner}})} \cdot \mathcal{O}_b^{(T_{\text{outer}})}$.
\end{theorem}

\begin{proof}
The inner phase $a$ contributes the phase angle $2\pi a/T_{\text{inner}}$ and the outer phase $b$ contributes the phase angle $2\pi b/T_{\text{outer}}$; their superposition gives the composite phase $\theta_{a,b} = 2\pi a/T_{\text{inner}} + 2\pi b/T_{\text{outer}}$, so the composite phase operator is
\[
\mathcal{O}_{a,b} = e^{i\theta_{a,b}} = e^{i2\pi a/T_{\text{inner}}} \cdot e^{i2\pi b/T_{\text{outer}}} = \mathcal{O}_a^{(T_{\text{inner}})} \cdot \mathcal{O}_b^{(T_{\text{outer}})}.
\]
Number of independent phases: the inner and outer phase indices form the Cartesian product $\mathbb{Z}_{T_{\text{inner}}} \times \mathbb{Z}_{T_{\text{outer}}}$, whose cardinality is $T_{\text{inner}} \times T_{\text{outer}}$, corresponding to the $T_{\text{inner}} \times T_{\text{outer}}$ independent components of the tensor-product representation $\mathcal{P}_{T_{\text{inner}}} \otimes \mathcal{P}_{T_{\text{outer}}}$. When $\gcd(T_{\text{inner}}, T_{\text{outer}}) = 1$, the phase-sum projection retains all the independent phases (see note \ref{ch20_epicycle_gcd_boundary}).
\end{proof}

\begin{remark}[The phase sum of nested epicycles and the tensor-product boundary]\label{ch20_epicycle_gcd_boundary}
The phase-sum projection $e^{i(\theta_a + \theta_b)}$ maps the inner and outer phase points to the same unit circle, and distinct pairs $(a,b)$ may then coincide: the set of all values of the phase sum forms the cyclic subgroup $\langle e^{i2\pi/\operatorname{lcm}(T_{\text{inner}},T_{\text{outer}})}\rangle$, with $\operatorname{lcm}(T_{\text{inner}}, T_{\text{outer}})$ distinct points. If and only if $\gcd(T_{\text{inner}}, T_{\text{outer}}) = 1$ is $\operatorname{lcm} = T_{\text{inner}} \times T_{\text{outer}}$, and only then does the phase-sum projection retain all $T_{\text{inner}} \times T_{\text{outer}}$ independent phases. Yin-Yang $\times$ the Five Phases ($2 \times 5$) and the Four Symbols $\times$ \textit{Sancai} ($4 \times 3$) are coprime, so the phase sum immediately gives the Heavenly Stems $10$ and the Earthly Branches $12$; for the Eight Trigrams $\times$ the Eight Trigrams, since $\gcd(8,8) = 8$, the phase sum has only $8$ points, and the tensor-product (direct-product) structure $C_8 \times C_8$ is required in order to obtain the $64$ hexagrams.
\end{remark}

\begin{example}[Epicycle generation of the Heavenly Stems]\label{ex:tiangan_generation}
The Five-Phase epicycle is the outer circle ($T_{\text{outer}} = 5$) and the Yin-Yang epicycle is the inner circle ($T_{\text{inner}} = 2$). For every single step forward of the Five Phases ($72^\circ$), Yin-Yang rotates through one full turn ($360^\circ$). The number of independent phases is $2 \times 5 = 10$, exactly the period of the Heavenly Stems.
\end{example}

\begin{example}[Epicycle generation of the Earthly Branches]\label{ex:dizhi_generation}
The Four-Symbols epicycle is the outer circle ($T_{\text{outer}} = 4$) and the Sancai epicycle is the inner circle ($T_{\text{inner}} = 3$). The number of independent phases is $3 \times 4 = 12$, exactly the period of the Earthly Branches.
\end{example}

\begin{example}[Epicycle generation of the Sixty-Four Hexagrams]\label{ex:liushisigua_generation}
The Eight-Trigrams epicycle is the outer circle ($T_{\text{outer}} = 8$) and the Eight-Trigrams epicycle is the inner circle ($T_{\text{inner}} = 8$). The number of independent phases is $8 \times 8 = 64$. The group structure is $C_8 \times C_8$.
\end{example}

\begin{lemma}[Parity filtering of the sexagenary cycle]\label{ch20_ganzhi_parity_filter}
Let $n \in \mathbb{Z}_{10}$ be the Heavenly-Stem index and $m \in \mathbb{Z}_{12}$ the Earthly-Branch index, with parity maps $\varphi_2: \mathbb{Z}_{10} \to \mathbb{Z}_2,\; n \mapsto n \bmod 2$ and $\psi_2: \mathbb{Z}_{12} \to \mathbb{Z}_2,\; m \mapsto m \bmod 2$. Then there are exactly $60$ effective combinations satisfying the synchronisation condition $n \equiv m \pmod 2$, and they correspond one-to-one with the sexagenary cycle.
\end{lemma}

\begin{proof}
The set of effective combinations is the fibre product
\[
E = \{(n,m) \in \mathbb{Z}_{10} \times \mathbb{Z}_{12} : \varphi_2(n) = \psi_2(m)\}.
\]
It is the kernel of the homomorphism $\Delta: \mathbb{Z}_{10} \times \mathbb{Z}_{12} \to \mathbb{Z}_2,\; (n,m) \mapsto \varphi_2(n) - \psi_2(m)$. Since both $10$ and $12$ contain the factor $2$, the maps $\varphi_2$ and $\psi_2$ are surjective, hence so is $\Delta$, and therefore $|E| = |\mathbb{Z}_{10} \times \mathbb{Z}_{12}| / |\mathbb{Z}_2| = 120/2 = 60$. A direct count confirms this: $\mathbb{Z}_{10}$ contains $5$ even and $5$ odd elements, and $\mathbb{Z}_{12}$ contains $6$ even and $6$ odd elements, so $|E| = 5\cdot 6 + 5\cdot 6 = 60$.

The diagonal embedding $\iota: \mathbb{Z}_{60} \to E,\; k \mapsto (k \bmod 10, k \bmod 12)$ is well defined (the parity of $k$ is the same under the two reductions) and injective: if $k \equiv k' \pmod{10}$ and $k \equiv k' \pmod{12}$, then $k \equiv k' \pmod{\operatorname{lcm}(10,12) = 60}$. Since $|E| = 60 = |\mathbb{Z}_{60}|$, the map $\iota$ is bijective, so $E \cong C_{60}$. Term by term in $k$ (the Heavenly-Stem index is counted from Jia and the Earthly-Branch index from Yin, so that $k$ is the index of the sexagenary cycle and Jiazi has index $10$): $k=0 \mapsto (0,0)$ (Jiayin, index $0$), $k=1 \mapsto (1,1)$ (Yimao), \dots, $k=10 \mapsto (0,10)$ (Jiazi, index $10$), which is precisely the sexagenary sequence beginning with Jiayin.
\end{proof}

\section{Prime factorisation and tensor structure}

\begin{theorem}[Direct-product decomposition of a cyclic group]\label{thm:cyclic_direct_product_yijing}
Let $T = \prod_{j=1}^r m_j$ with $m_1, \dots, m_r$ pairwise coprime. Then $C_T \cong C_{m_1} \times \cdots \times C_{m_r}$.
\end{theorem}

\begin{proof}
By the Chinese remainder theorem: when $m_1, \dots, m_r$ are pairwise coprime, the map
\[
\varphi: \mathbb{Z}_T \to \mathbb{Z}_{m_1} \times \cdots \times \mathbb{Z}_{m_r}, \quad k \mapsto (k \bmod m_1, \dots, k \bmod m_r)
\]
is a ring isomorphism, and in particular a group isomorphism. Hence $C_T \cong C_{m_1} \times \cdots \times C_{m_r}$.
\end{proof}

\begin{definition}[Tensor representation of a phase operator]\label{def:tensor_representation_yijing}
Under the direct-product decomposition, every phase operator is represented as an $r$-th order tensor:
\[
\mathbf{O}_{a_1,\dots,a_r} = \exp\left(i2\pi \sum_{j=1}^r \frac{a_j}{m_j}\right).
\]
\end{definition}

\begin{proposition}[Tensor product of irreducible representations]\label{prop:irrep_tensor_product_yijing}
Let $G = G_1 \times \cdots \times G_r$ be a direct product of cyclic group factors (the group structures of the thirteen symbol systems of this chapter are all of this kind). Then every irreducible representation of $G$ is the tensor product of irreducible representations of the factors: $\rho = \rho_1 \otimes \cdots \otimes \rho_r$, and $\widehat{G} \cong \widehat{G_1} \times \cdots \times \widehat{G_r}$.
\end{proposition}

\begin{proof}
By theorem \ref{thm:cyclic_direct_product_yijing}, when $T = \prod_j m_j$ with the factors pairwise coprime, $C_T \cong C_{m_1} \times \cdots \times C_{m_r}$. Let the generators of the factors be $g_j$ ($g_j^{m_j} = 1$); then every element of $G$ is uniquely written as $g_1^{x_1} \cdots g_r^{x_r}$. Any one-dimensional irreducible representation $\rho$ of $G$ is uniquely determined by $\rho(g_1),\dots,\rho(g_r)$; since $\rho(g_j)^{m_j} = \rho(g_j^{m_j}) = \rho(1) = 1$, the number $\rho(g_j)$ is an $m_j$-th root of unity, so there exists $a_j \in \mathbb{Z}_{m_j}$ with $\rho(g_j) = e^{i2\pi a_j/m_j}$. Hence
\[
\rho(g_1^{x_1}\cdots g_r^{x_r}) = \prod_{j=1}^r e^{i2\pi a_j x_j / m_j},
\]
which is precisely the value of the tensor product of the irreducible representations $\rho_j(g_j^{x_j}) = e^{i2\pi a_j x_j/m_j}$ of the factors. The indices $(a_1,\dots,a_r)$ run independently through $\mathbb{Z}_{m_1}\times\cdots\times\mathbb{Z}_{m_r}$, giving $\prod_j m_j = T$ combinations in all, equal to the total number $T$ of irreducible representations of $C_T$; hence there is neither omission nor repetition.
\end{proof}

\section{He Tu complex functions}

This section establishes the \textbf{He Tu complex function}---a complex function of period $10$. The He Tu structure reveals that the Fourier decomposition of period $10$ automatically exhibits the algebraic grading $\mathbb{Z}_{10} \cong \mathbb{Z}_2 \times \mathbb{Z}_5$, in which the $\mathbb{Z}_5$ component corresponds to the mutual-generation cycle of the \textbf{Five Phases} and the $\mathbb{Z}_2$ component to the twofold alternation of \textbf{Yin-Yang}.

\begin{definition}[He Tu number--Five-Phase correspondence]\label{def:hetu_mapping_yijing}
Define the \textbf{He Tu map} $\varphi: \mathbb{Z}_{10} \to \mathbb{Z}_5$ by
\[
\varphi(n) := n \bmod 5, \qquad n \in \mathbb{Z}_{10} = \{0, 1, \dots, 9\},
\]
where the remainder $0$ is by convention associated with `earth' among the Five Phases. Specifically:
\[
\begin{array}{c|ccccc}
\varphi(n) & 0 & 1 & 2 & 3 & 4 \\
\hline
\text{Five Phases} & \text{earth} & \text{water} & \text{fire} & \text{wood} & \text{metal} \\
\text{Numbers} & 5, 10 & 1, 6 & 2, 7 & 3, 8 & 4, 9
\end{array}
\]
Moreover, the \textbf{Yin-Yang} component is given by the parity of $n$: an odd $n$ corresponds to yang and an even $n$ to yin. Thus the He Tu map is a group homomorphism
\[
\varphi: \mathbb{Z}_{10} \to \mathbb{Z}_5, \qquad \ker \varphi = \{0, 5\} \cong \mathbb{Z}_2.
\]
\end{definition}

\begin{remark}[He Tu complex functions (the same definition as in the `Series and Convolution' chapter)]\label{def:hetu_function_yijing}
The He Tu complex function---that is, a complex function satisfying $f(z + 10) = f(z)$---is the \textbf{same definition} as \cref{def:hetu_function} of the `Series and Convolution' chapter, and its Fourier expansion $f(z) = \sum_{n} c_n e^{2\pi i n z/10}$ and the coefficient formula are likewise exactly the same, so they are not repeated here.
\end{remark}

\begin{remark}[The He Tu Fourier decomposition (restatement)]\label{thm:hetu_fourier_decomposition_yijing}
This theorem is the same conclusion as \cref{thm:hetu_fourier_decomposition} of the `Series and Convolution' chapter: a uniformly convergent He Tu complex function can be uniquely decomposed as
\[
f(z) = \sum_{k=0}^{4} \Bigl( c_k + c_{k+5}\, e^{\pi i z} \Bigr) e^{2\pi i k z / 10},
\]
where $e^{\pi i z}$ is the Yin-Yang factor of period $2$. Both the statement and the proof are to be found there; the reader may cite the conclusion of this section directly from that place.
\end{remark}

\begin{proof}
The proof of this theorem is exactly the same as the proof of \cref{thm:hetu_fourier_decomposition} of the `Series and Convolution' chapter (where the substitution computation for $\mathbb{Z}_{10} \cong \mathbb{Z}_2 \times \mathbb{Z}_5$ is given term by term), and is not repeated here.
\end{proof}

\begin{definition}[Five-Phase components]\label{def:hetu_five_components_yijing}
In the above decomposition,
\[
f_k(z) := \Bigl( c_k + c_{k+5}\, e^{\pi i z} \Bigr)\, e^{2\pi i k z / 10}, \qquad k = 0, 1, 2, 3, 4
\]
is called the \textbf{$k$-th Five-Phase component} of $f$. By the He Tu map (definition \ref{def:hetu_mapping_yijing}), the five components correspond in order to:
\[
f_0 \leftrightarrow \text{earth},\quad f_1 \leftrightarrow \text{water},\quad f_2 \leftrightarrow \text{fire},\quad f_3 \leftrightarrow \text{wood},\quad f_4 \leftrightarrow \text{metal}.
\]
Each Five-Phase component consists of two parts: $c_k\, e^{2\pi i k z / 10}$ (the yang component) and $c_{k+5}\, e^{\pi i z} e^{2\pi i k z / 10}$ (the yin component).
\end{definition}

\begin{remark}[Period structure of the Five-Phase components (restated in five-phase notation)]\label{prop:hetu_five_periods_yijing}
This proposition is a restatement of \cref{prop:hetu_five_periods} of the `Series and Convolution' chapter in \textbf{five-phase notation}, with exactly the same mathematical content; what is given after the proof there is the period table indexed by $k$, and here only the index is replaced by the names of the Five Phases:
\[
\begin{array}{c|ccccc}
k & 0 & 1 & 2 & 3 & 4 \\
\hline
\text{Five Phases} & \text{earth} & \text{water} & \text{fire} & \text{wood} & \text{metal} \\
\text{fundamental period} & 1 & 10 & 5 & 10 & 5
\end{array}
\]
That is: $f_0$ (earth) contains only the constant term; $f_1, f_3$ (water, wood) are components of period $10$; $f_2, f_4$ (fire, metal) are components of period $5$. The proof is to be found in \cref{prop:hetu_five_periods} and is not repeated here.
\end{remark}

\begin{remark}[The yin--yang flip of a He Tu complex function (restated in five-phase notation)]\label{prop:hetu_yin_yang_flip_yijing}
This proposition is the same proposition as \cref{prop:hetu_yin_yang_flip} of the `Series and Convolution' chapter, and its statement and proof are not repeated here; in five-phase notation it reads as follows: the half-period translation $z \mapsto z + 5$ gives each component a factor $(-1)^k$; the even components ($k = 0, 2, 4$, corresponding to earth, fire and metal) are unchanged, whilst the odd components ($k = 1, 3$, corresponding to water and wood) change sign as a whole.
\end{remark}

\begin{remark}[The He Tu phase operator (the same definition as in the `Series and Convolution' chapter)]\label{def:hetu_phase_operator_yijing}
The He Tu phase operator $\mathcal{O}_n^{(10)} = e^{2\pi i n / 10}$ is the \textbf{same object} as \cref{def:hetu_phase_operator} of the `Series and Convolution' chapter, and its decomposition $\mathbb{Z}_{10} \cong \mathbb{Z}_2 \times \mathbb{Z}_5$ is likewise exactly the same, so it is not defined again here.

The I Ching context merely gives these two factors other names: the $\mathbb{Z}_2$ component $(-1)^{\varepsilon}$ is called the \textbf{Yin-Yang sign}, and the $\mathbb{Z}_5$ component $e^{2\pi i k / 10}$ is called the \textbf{Five-Phase phase}.
\end{remark}

\begin{remark}[Direct-product decomposition of the He Tu phase-operator group (restatement)]\label{prop:hetu_phase_group_yijing}
This proposition is the same proposition as \cref{prop:hetu_phase_group} of the `Series and Convolution' chapter: the phase-operator group $\{\mathcal{O}_n^{(10)}\}_{n=0}^{9}$ is isomorphic to $C_{10} \cong C_2 \times C_5$, where $C_2$ corresponds to Yin-Yang and $C_5$ to the Five Phases. Both the statement and the proof are to be found there and are not repeated here.
\end{remark}

\begin{remark}[DFT diagonalisation of the He Tu circulant matrix (restatement)]\label{thm:hetu_circulant_yijing}
This theorem is the same conclusion as \cref{thm:hetu_circulant} of the `Series and Convolution' chapter: the $10\times 10$ circulant matrix with first-row generator $(h_0,\dots,h_9)$ is unitarily diagonalised by the DFT matrix, and the eigenvalues $\hat h_n$ are divided into a yin group ($\varepsilon=1$) and a yang group ($\varepsilon=0$) according to $n = 5\varepsilon + k$. Both the statement and the proof are to be found there.
\end{remark}

\begin{proof}
The proof of this theorem is exactly the same as the proof of \cref{thm:hetu_circulant} of the `Series and Convolution' chapter (the DFT diagonalisation of the circulant matrix and the grading of the eigenvalues by $n = 5\varepsilon + k$), and is not repeated here.
\end{proof}

\begin{example}[Concrete decomposition of a He Tu complex function]\label{ex:hetu_concrete_yijing}
This example has the same content as Chapter~\ref{chap:series_conv}, the section `He Tu complex functions: the $\mathbb{Z}_2\times\mathbb{Z}_5$ decomposition of period $10$' (\cref{sec:hetu_function}), \cref{ex:hetu_concrete}; to avoid repeating the statement, only the reference anchor is retained here, and the complete statement is given there. The only difference between the two places is the column names of the component table (that place uses `label/component', this place uses `Five Phases/Yin-Yang'); the correspondence between the Five Phases and Yin-Yang is given by \cref{def:hetu_mapping}, and no new example is set up.
\end{example}

\begin{remark}[The algebraic meaning of the He Tu structure]
This note has the same content as Chapter~\ref{chap:series_conv}, the section `He Tu complex functions: the $\mathbb{Z}_2\times\mathbb{Z}_5$ decomposition of period $10$' (\cref{sec:hetu_function}), \cref{rem:hetu_algebraic_meaning}; to avoid repeating the statement, only the reference anchor is retained here, and the complete statement is given there. The I Ching context merely gives the two factors other names: the $\mathbb{Z}_2$ component $(-1)^{\varepsilon}$ is called the \textbf{Yin-Yang sign}, and the $\mathbb{Z}_5$ component $e^{2\pi i k/10}$ is called the \textbf{Five-Phase phase}.
\end{remark}

\section{The Lo Shu matrix}

This section uses all the conclusions concerning the \textbf{Lo Shu matrix} in the context of the *I Ching* symbol system. The definition of this matrix, its magic-square property, the characterisation of its balance point, its antipodal balance, its spectral structure and its translation--rotation decomposition have already been given in full and proved one by one. The items of this section are all citations of the conclusions of that section; their statements and proofs are not repeated, and only the numbering anchors needed for references within this chapter are retained; wherever `see that place' occurs below, it refers uniformly to Chapter~\ref{chap:linear_algebra}, \cref{sec:luoshu_matrix}.

\begin{definition}[The Lo Shu matrix]\label{def:luoshu_matrix_yijing}
This definition has the same content as Chapter~\ref{chap:linear_algebra}, the section `The Lo Shu matrix: a $3\times 3$ magic square with $5$ as the balance point' (\cref{sec:luoshu_matrix}), \cref{def:luoshu_matrix}; to avoid repeating the statement, only the reference anchor is retained here, and the complete statement is given there.
\end{definition}

\begin{definition}[Row sums, column sums and diagonal sums]\label{def:luoshu_line_sums_yijing}
This definition has the same content as Chapter~\ref{chap:linear_algebra}, the section `The Lo Shu matrix: a $3\times 3$ magic square with $5$ as the balance point' (\cref{sec:luoshu_matrix}), \cref{def:luoshu_line_sums}; to avoid repeating the statement, only the reference anchor is retained here, and the complete statement is given there.
\end{definition}

\begin{proposition}[The magic-square property of the Lo Shu matrix]\label{prop:luoshu_magic_sum_yijing}
This proposition has the same content as Chapter~\ref{chap:linear_algebra}, the section `The Lo Shu matrix: a $3\times 3$ magic square with $5$ as the balance point' (\cref{sec:luoshu_matrix}), \cref{prop:luoshu_magic_sum}; to avoid repeating the statement, only the reference anchor is retained here, and the complete statement and proof are given there.
\end{proposition}

\begin{definition}[The Lo Shu balance point]\label{def:luoshu_balance_point_yijing}
This definition has the same content as Chapter~\ref{chap:linear_algebra}, the section `The Lo Shu matrix: a $3\times 3$ magic square with $5$ as the balance point' (\cref{sec:luoshu_matrix}), \cref{def:luoshu_balance_point}; to avoid repeating the statement, only the reference anchor is retained here, and the complete statement is given there.
\end{definition}

\begin{proposition}[Threefold characterisation of the balance point]\label{prop:luoshu_balance_triple_yijing}
This proposition has the same content as Chapter~\ref{chap:linear_algebra}, the section `The Lo Shu matrix: a $3\times 3$ magic square with $5$ as the balance point' (\cref{sec:luoshu_matrix}), \cref{prop:luoshu_balance_triple}; to avoid repeating the statement, only the reference anchor is retained here, and the complete statement and proof are given there.
\end{proposition}

\begin{proposition}[Antipodal balance of the Lo Shu matrix]\label{prop:luoshu_antipodal_yijing}
This proposition has the same content as Chapter~\ref{chap:linear_algebra}, the section `The Lo Shu matrix: a $3\times 3$ magic square with $5$ as the balance point' (\cref{sec:luoshu_matrix}), \cref{prop:luoshu_antipodal}; to avoid repeating the statement, only the reference anchor is retained here, and the complete statement and proof are given there.
\end{proposition}

\begin{theorem}[Eigenvalues and spectral structure of the Lo Shu matrix]\label{thm:luoshu_eigenvalues_yijing}
This theorem has the same content as Chapter~\ref{chap:linear_algebra}, the section `The Lo Shu matrix: a $3\times 3$ magic square with $5$ as the balance point' (\cref{sec:luoshu_matrix}), \cref{thm:luoshu_eigenvalues}; to avoid repeating the statement, only the reference anchor is retained here, and the complete statement and proof are given there.
\end{theorem}

\begin{theorem}[Translation--rotation decomposition of the Lo Shu matrix]\label{thm:luoshu_decomposition_yijing}
This theorem has the same content as Chapter~\ref{chap:linear_algebra}, the section `The Lo Shu matrix: a $3\times 3$ magic square with $5$ as the balance point' (\cref{sec:luoshu_matrix}), \cref{thm:luoshu_decomposition}; to avoid repeating the statement, only the reference anchor is retained here, and the complete statement and proof are given there.
\end{theorem}

\begin{corollary}[Structural meaning of the Lo Shu balance point and the magic sum]\label{cor:luoshu_structure}
This corollary has the same content as Chapter~\ref{chap:linear_algebra}, the section `The Lo Shu matrix: a $3\times 3$ magic square with $5$ as the balance point' (\cref{sec:luoshu_matrix}), \cref{cor:luoshu_grammar}; to avoid repeating the statement, only the reference anchor is retained here, and the complete statement and proof are given there.
\end{corollary}

\begin{example}[Spectral decomposition of the Lo Shu matrix]\label{ex:luoshu_spectral_yijing}
This example has the same content as Chapter~\ref{chap:linear_algebra}, the section `The Lo Shu matrix: a $3\times 3$ magic square with $5$ as the balance point' (\cref{sec:luoshu_matrix}), \cref{ex:luoshu_spectral}; to avoid repeating the statement, only the reference anchor is retained here, and the complete statement is given there.
\end{example}

\begin{remark}[Interface between the Lo Shu matrix and the He Tu structure]\label{rem:luoshu_hetu_interface_yijing}
The factor $5$ in the magic sum $15 = 3 \times 5$ of the Lo Shu matrix is algebraically related to the $\mathbb{Z}_5$ component of $\mathbb{Z}_{10} \cong \mathbb{Z}_2 \times \mathbb{Z}_5$ in the He Tu structure: the Lo Shu centre $5$ corresponds to the apex of the Five-Phase cycle, and the antipodal balance of the Lo Shu matrix (with sum $10 = 2 \times 5$) corresponds to the twofold alternation of Yin-Yang. The purely rotational eigenvalues $\pm 2\sqrt{6}i$ of the Lo Shu matrix on $\mathbf{1}^\perp$ provide a concrete example of a $3$-dimensional rotation, whose period $\pi/\sqrt{6}$ is completely determined by the algebraic structure of the matrix.
\end{remark}

\section{Table of phase operators of the thirteen symbol systems}

\subsection{\texorpdfstring{Taiji ($T=1$)}{Taiji (T=1)}}

\begin{definition}[Taiji phase operators]\label{def:taiji_operators}
The Taiji period is $T=1$, the bit width is $W=1$ and the bandwidth is $\Delta\theta=2\pi$. The group structure is $C_1$ and the tensor shape is $1$. There is only one phase point, which describes `existence'.
\end{definition}

\begin{center}
\begin{tabular}{>{\centering\arraybackslash}p{1.2cm}>{\centering\arraybackslash}p{2cm}>{\centering\arraybackslash}p{2.5cm}>{\centering\arraybackslash}p{3.5cm}}
\toprule
$k$ & Symbol & Phase $\theta_k$ & Operator $\mathcal{O}_k^{(1)}$ \\
\midrule
$0$ & Taiji & $0$ & $1+0i$ \\
\bottomrule
\end{tabular}
\end{center}

\subsection{\texorpdfstring{The Two Modes ($T=2$)}{The Two Modes (T=2)}}

\begin{definition}[Yin-Yang phase operators]\label{def:yinyang_operators}
The Yin-Yang period is $T=2$, the bit width is $W=2$ and the bandwidth is $\Delta\theta=\pi$. The group structure is $C_2$ and the tensor shape is $2$. The all-yin line configuration ($k=0$) is Zhen (the beginning of wood) and the all-yang line configuration ($k=1$) is Kan (the end of water).
\end{definition}

\begin{center}
\begin{tabular}{>{\centering\arraybackslash}p{1.2cm}>{\centering\arraybackslash}p{1.2cm}>{\centering\arraybackslash}p{1.5cm}>{\centering\arraybackslash}p{2cm}>{\centering\arraybackslash}p{2.5cm}>{\centering\arraybackslash}p{3cm}}
\toprule
$k$ & Binary & Symbol & Five Phases & Phase $\theta_k$ & Operator $\mathcal{O}_k^{(2)}$ \\
\midrule
$0$ & $0$ & yin & Zhen (beginning of wood) & $0$ & $1+0i$ \\
$1$ & $1$ & yang & Kan (end of water) & $\pi$ & $-1+0i$ \\
\bottomrule
\end{tabular}
\end{center}

\subsection{\texorpdfstring{Sancai ($T=3$)}{Sancai (T=3)}}

\begin{definition}[Sancai phase operators]\label{def:sancai_operators}
The Sancai period is $T=3$, the bit width is $W=3$ and the bandwidth is $\Delta\theta=2\pi/3$. The group structure is $C_3$ and the tensor shape is $3$.
\end{definition}

\begin{center}
\begin{tabular}{>{\centering\arraybackslash}p{1.2cm}>{\centering\arraybackslash}p{1.5cm}>{\centering\arraybackslash}p{1.5cm}>{\centering\arraybackslash}p{2.5cm}>{\centering\arraybackslash}p{4cm}}
\toprule
$k$ & Symbol & Five Phases & Phase $\theta_k$ & Operator $\mathcal{O}_k^{(3)}$ \\
\midrule
$0$ & heaven & yang wood (beginning) & $0$ & $1+0i$ \\
$1$ & human & fire & $2\pi/3$ & $-0.5+0.8660i$ \\
$2$ & earth & yin water (end) & $4\pi/3$ & $-0.5-0.8660i$ \\
\bottomrule
\end{tabular}
\end{center}

\subsection{\texorpdfstring{The Four Symbols ($T=4$)}{The Four Symbols (T=4)}}

\begin{definition}[Four-Symbols phase operators]\label{def:sixiang_operators}
The Four-Symbols period is $T=4=2^2$, the bit width is $W=4$ and the bandwidth is $\Delta\theta=\pi/2$. The group structure is $C_4$ and the tensor shape is $4$.
\end{definition}

\begin{center}
\begin{tabular}{>{\centering\arraybackslash}p{1.2cm}>{\centering\arraybackslash}p{1.2cm}>{\centering\arraybackslash}p{1.5cm}>{\centering\arraybackslash}p{1.5cm}>{\centering\arraybackslash}p{2.5cm}>{\centering\arraybackslash}p{3cm}}
\toprule
$k$ & Binary & Symbol & Five Phases & Phase $\theta_k$ & Operator $\mathcal{O}_k^{(4)}$ \\
\midrule
$0$ & $00$ & greater yin & Zhen (beginning of wood) & $0$ & $1+0i$ \\
$1$ & $01$ & lesser yang & wood & $\pi/2$ & $0+1i$ \\
$2$ & $10$ & lesser yin & metal & $\pi$ & $-1+0i$ \\
$3$ & $11$ & greater yang & Kan (end of water) & $3\pi/2$ & $0-1i$ \\
\bottomrule
\end{tabular}
\end{center}

\subsection{\texorpdfstring{The Five Phases ($T=5$)}{The Five Phases (T=5)}}

\begin{definition}[Five-Phase phase operators]\label{def:wuxing_operators}
The Five-Phase period is $T=5$, the bit width is $W=5$ and the bandwidth is $\Delta\theta=2\pi/5$. The group structure is $C_5$ and the tensor shape is $5$.
\end{definition}

\begin{center}
\begin{tabular}{>{\centering\arraybackslash}p{1.2cm}>{\centering\arraybackslash}p{1.5cm}>{\centering\arraybackslash}p{1.5cm}>{\centering\arraybackslash}p{2.5cm}>{\centering\arraybackslash}p{4cm}}
\toprule
$k$ & Symbol & Five Phases & Phase $\theta_k$ & Operator $\mathcal{O}_k^{(5)}$ \\
\midrule
$0$ & wood & yang wood (beginning) & $0$ & $1+0i$ \\
$1$ & fire & fire & $2\pi/5$ & $0.3090+0.9511i$ \\
$2$ & earth & earth & $4\pi/5$ & $-0.8090+0.5878i$ \\
$3$ & metal & metal & $6\pi/5$ & $-0.8090-0.5878i$ \\
$4$ & water & yin water (end) & $8\pi/5$ & $0.3090-0.9511i$ \\
\bottomrule
\end{tabular}
\end{center}

\subsection{\texorpdfstring{The Seven Luminaries ($T=7$)}{The Seven Luminaries (T=7)}}

\begin{definition}[Seven-Luminaries phase operators]\label{def:qizheng_operators}
The Seven-Luminaries period is $T=7$, the bit width is $W=7$ and the bandwidth is $\Delta\theta=2\pi/7$. The group structure is $C_7$ and the tensor shape is $7$.
\end{definition}

\begin{center}
\begin{tabular}{>{\centering\arraybackslash}p{1.2cm}>{\centering\arraybackslash}p{1.5cm}>{\centering\arraybackslash}p{1.5cm}>{\centering\arraybackslash}p{2.5cm}>{\centering\arraybackslash}p{4cm}}
\toprule
$k$ & Symbol & Five Phases & Phase $\theta_k$ & Operator $\mathcal{O}_k^{(7)}$ \\
\midrule
$0$ & Sun & yang wood (beginning) & $0$ & $1+0i$ \\
$1$ & Moon & --- & $2\pi/7$ & $0.6235+0.7818i$ \\
$2$ & Mars & fire & $4\pi/7$ & $-0.2225+0.9749i$ \\
$3$ & Mercury & water & $6\pi/7$ & $-0.9010+0.4339i$ \\
$4$ & Jupiter & wood & $8\pi/7$ & $-0.9010-0.4339i$ \\
$5$ & Venus & metal & $10\pi/7$ & $-0.2225-0.9749i$ \\
$6$ & Saturn & yin water (end) & $12\pi/7$ & $0.6235-0.7818i$ \\
\bottomrule
\end{tabular}
\end{center}

\subsection{\texorpdfstring{The Eight Trigrams ($T=8$)}{The Eight Trigrams (T=8)}}

\begin{definition}[Eight-Trigrams phase operators]\label{def:bagua_operators}
The Eight-Trigrams period is $T=8=2^3$, the bit width is $W=8$ and the bandwidth is $\Delta\theta=\pi/4$. The group structure is $C_8$ and the tensor shape is $8$. The all-yin line configuration ($000$, $k=0$, Zhen) is the beginning of wood, and the all-yang line configuration ($111$, $k=7$, Kan) is the end of water.
\end{definition}

\begin{center}
\begin{tabular}{>{\centering\arraybackslash}p{1cm}>{\centering\arraybackslash}p{1.2cm}>{\centering\arraybackslash}p{1cm}>{\centering\arraybackslash}p{1.2cm}>{\centering\arraybackslash}p{1.2cm}>{\centering\arraybackslash}p{2cm}>{\centering\arraybackslash}p{3.5cm}}
\toprule
$k$ & Binary & Line image & Hexagram name & Five Phases & Phase $\theta_k$ & Operator $\mathcal{O}_k^{(8)}$ \\
\midrule
$0$ & $000$ & ☳ & Zhen (Thunder) & wood (beginning) & $0$ & $1+0i$ \\
$1$ & $001$ & ☴ & Xun (Wind) & wood & $\pi/4$ & $0.7071+0.7071i$ \\
$2$ & $010$ & ☶ & Gen (Mountain) & earth & $\pi/2$ & $0+1i$ \\
$3$ & $011$ & ☲ & Li (Fire) & fire & $3\pi/4$ & $-0.7071+0.7071i$ \\
$4$ & $100$ & ☰ & Qian (Heaven) & metal & $\pi$ & $-1+0i$ \\
$5$ & $101$ & ☷ & Kun (Earth) & earth & $5\pi/4$ & $-0.7071-0.7071i$ \\
$6$ & $110$ & ☱ & Dui (Lake) & metal & $3\pi/2$ & $0-1i$ \\
$7$ & $111$ & ☵ & Kan (Water) & water (end) & $7\pi/4$ & $0.7071-0.7071i$ \\
\bottomrule
\end{tabular}
\end{center}

\subsection{\texorpdfstring{The Heavenly Stems ($T=10$)}{The Heavenly Stems (T=10)}}

\begin{definition}[Heavenly-Stems phase operators]\label{def:tiangan_operators}
The Heavenly-Stems period is $T=10=2 \times 5$, the bit width is $W=10$ and the bandwidth is $\Delta\theta=\pi/5$. The group structure is $C_{10} \cong C_2 \times C_5$ and the tensor shape is $2 \times 5$. It is generated by nested epicycles of Yin-Yang $\times$ the Five Phases.
\end{definition}

\begin{center}
\begin{tabular}{>{\centering\arraybackslash}p{1cm}>{\centering\arraybackslash}p{1.2cm}>{\centering\arraybackslash}p{1.8cm}>{\centering\arraybackslash}p{1.8cm}>{\centering\arraybackslash}p{2cm}>{\centering\arraybackslash}p{3.5cm}}
\toprule
$k$ & Symbol & Yin-Yang and Five Phases & Tensor index & Phase $\theta_k$ & Operator $\mathcal{O}_k^{(10)}$ \\
\midrule
$0$ & Jia & yang wood (beginning) & $(0,0)$ & $0$ & $1+0i$ \\
$1$ & Yi & yin wood & $(0,1)$ & $\pi/5$ & $0.8090+0.5878i$ \\
$2$ & Bing & yang fire & $(0,2)$ & $2\pi/5$ & $0.3090+0.9511i$ \\
$3$ & Ding & yin fire & $(0,3)$ & $3\pi/5$ & $-0.3090+0.9511i$ \\
$4$ & Wu & yang earth & $(0,4)$ & $4\pi/5$ & $-0.8090+0.5878i$ \\
$5$ & Ji & yin earth & $(1,0)$ & $\pi$ & $-1+0i$ \\
$6$ & Geng & yang metal & $(1,1)$ & $6\pi/5$ & $-0.8090-0.5878i$ \\
$7$ & Xin & yin metal & $(1,2)$ & $7\pi/5$ & $-0.3090-0.9511i$ \\
$8$ & Ren & yang water & $(1,3)$ & $8\pi/5$ & $0.3090-0.9511i$ \\
$9$ & Gui & yin water (end) & $(1,4)$ & $9\pi/5$ & $0.8090-0.5878i$ \\
\bottomrule
\end{tabular}
\end{center}

\subsection{\texorpdfstring{The Earthly Branches ($T=12$)}{The Earthly Branches (T=12)}}

\begin{definition}[Earthly-Branches phase operators]\label{def:dizhi_operators}
The Earthly-Branches period is $T=12=4 \times 3$, the bit width is $W=12$ and the bandwidth is $\Delta\theta=\pi/6$. The group structure is $C_{12} \cong C_4 \times C_3$ and the tensor shape is $4 \times 3$. It is generated by nested epicycles of the Four Symbols $\times$ Sancai.
\end{definition}

\begin{center}
\begin{tabular}{>{\centering\arraybackslash}p{1cm}>{\centering\arraybackslash}p{1.2cm}>{\centering\arraybackslash}p{1.5cm}>{\centering\arraybackslash}p{1.8cm}>{\centering\arraybackslash}p{2cm}>{\centering\arraybackslash}p{3.5cm}}
\toprule
$k$ & Symbol & Five Phases & Tensor index & Phase $\theta_k$ & Operator $\mathcal{O}_k^{(12)}$ \\
\midrule
$0$ & Yin & yang wood (beginning) & $(0,0)$ & $0$ & $1+0i$ \\
$1$ & Mao & yin wood & $(0,1)$ & $\pi/6$ & $0.8660+0.5i$ \\
$2$ & Chen & yang earth & $(0,2)$ & $\pi/3$ & $0.5+0.8660i$ \\
$3$ & Si & yin fire & $(1,0)$ & $\pi/2$ & $0+1i$ \\
$4$ & Wu & yang fire & $(1,1)$ & $2\pi/3$ & $-0.5+0.8660i$ \\
$5$ & Wei & yin earth & $(1,2)$ & $5\pi/6$ & $-0.8660+0.5i$ \\
$6$ & Shen & yang metal & $(2,0)$ & $\pi$ & $-1+0i$ \\
$7$ & You & yin metal & $(2,1)$ & $7\pi/6$ & $-0.8660-0.5i$ \\
$8$ & Xu & yang earth & $(2,2)$ & $4\pi/3$ & $-0.5-0.8660i$ \\
$9$ & Hai & yin water & $(3,0)$ & $3\pi/2$ & $0-1i$ \\
$10$ & Zi & yang water & $(3,1)$ & $5\pi/3$ & $0.5-0.8660i$ \\
$11$ & Chou & yin water (end) & $(3,2)$ & $11\pi/6$ & $0.8660-0.5i$ \\
\bottomrule
\end{tabular}
\end{center}

\subsection{\texorpdfstring{The twenty-four solar terms ($T=24$)}{The twenty-four solar terms (T=24)}}

\begin{definition}[Phase operators of the twenty-four solar terms]\label{def:jieqi_operators}
The period of the twenty-four solar terms is $T=24=8 \times 3$, the bit width is $W=24$ and the bandwidth is $\Delta\theta=\pi/12$. The group structure is $C_{24} \cong C_8 \times C_3$ and the tensor shape is $8 \times 3$. They are generated by nested epicycles of the Eight Trigrams $\times$ Sancai.
\end{definition}

\begin{center}
\begin{tabular}{>{\centering\arraybackslash}p{1cm}>{\centering\arraybackslash}p{1.5cm}>{\centering\arraybackslash}p{1.5cm}>{\centering\arraybackslash}p{2cm}>{\centering\arraybackslash}p{3.5cm}}
\toprule
$k$ & Symbol & Five Phases & Phase $\theta_k$ & Operator $\mathcal{O}_k^{(24)}$ \\
\midrule
$0$ & Lichun & yang wood (beginning) & $0$ & $1+0i$ \\
$1$ & Yushui & yin wood & $\pi/12$ & $0.9659+0.2588i$ \\
$2$ & Jingzhe & yang wood & $\pi/6$ & $0.8660+0.5i$ \\
$3$ & Chunfen & yin wood & $\pi/4$ & $0.7071+0.7071i$ \\
$4$ & Qingming & yang wood & $\pi/3$ & $0.5+0.8660i$ \\
$5$ & Guyu & yin wood & $5\pi/12$ & $0.2588+0.9659i$ \\
$6$ & Lixia & yang fire & $\pi/2$ & $0+1i$ \\
$7$ & Xiaoman & yin fire & $7\pi/12$ & $-0.2588+0.9659i$ \\
$8$ & Mangzhong & yang fire & $2\pi/3$ & $-0.5+0.8660i$ \\
$9$ & Xiazhi & yin fire & $3\pi/4$ & $-0.7071+0.7071i$ \\
$10$ & Xiaoshu & yang fire & $5\pi/6$ & $-0.8660+0.5i$ \\
$11$ & Dashu & yin fire & $11\pi/12$ & $-0.9659+0.2588i$ \\
$12$ & Liqiu & yang metal & $\pi$ & $-1+0i$ \\
$13$ & Chushu & yin metal & $13\pi/12$ & $-0.9659-0.2588i$ \\
$14$ & Bailu & yang metal & $7\pi/6$ & $-0.8660-0.5i$ \\
$15$ & Qiufen & yin metal & $5\pi/4$ & $-0.7071-0.7071i$ \\
$16$ & Hanlu & yang metal & $4\pi/3$ & $-0.5-0.8660i$ \\
$17$ & Shuangjiang & yin metal & $17\pi/12$ & $-0.2588-0.9659i$ \\
$18$ & Lidong & yang water & $3\pi/2$ & $0-1i$ \\
$19$ & Xiaoxue & yin water & $19\pi/12$ & $0.2588-0.9659i$ \\
$20$ & Daxue & yang water & $5\pi/3$ & $0.5-0.8660i$ \\
$21$ & Dongzhi & yin water & $7\pi/4$ & $0.7071-0.7071i$ \\
$22$ & Xiaohan & yang water & $11\pi/6$ & $0.8660-0.5i$ \\
$23$ & Dahan & yin water (end) & $23\pi/12$ & $0.9659-0.2588i$ \\
\bottomrule
\end{tabular}
\end{center}

\subsection{\texorpdfstring{The thirty Na-Yin ($T=30$)}{The thirty Na-Yin (T=30)}}

\begin{definition}[Thirty-Na-Yin phase operators]\label{def:nayin_operators}
The period of the thirty Na-Yin is $T=30=2 \times 3 \times 5$, the bit width is $W=30$ and the bandwidth is $\Delta\theta=\pi/15$. The group structure is $C_{30} \cong C_2 \times C_3 \times C_5$ and the tensor shape is $2 \times 3 \times 5$. They are generated by nested epicycles of Yin-Yang $\times$ Sancai $\times$ the Five Phases.
\end{definition}

\begin{center}
\begin{tabular}{>{\centering\arraybackslash}p{1cm}>{\centering\arraybackslash}p{2.5cm}>{\centering\arraybackslash}p{3cm}>{\centering\arraybackslash}p{2cm}>{\centering\arraybackslash}p{3.5cm}}
\toprule
$k$ & Na-Yin & Stems and Branches & Phase $\theta_k$ & Operator $\mathcal{O}_k^{(30)}$ \\
\midrule
$0$  & Great Stream Water   & Jiayin, Yimao & $0$ & $1+0i$ \\
$1$  & Earth in the Sand   & Bingchen, Dingsi & $\pi/15$ & $0.9781+0.2079i$ \\
$2$  & Fire in the Sky   & Wuwu, Jiwei & $2\pi/15$ & $0.9135+0.4067i$ \\
$3$  & Pomegranate Wood   & Gengshen, Xinyou & $\pi/5$ & $0.8090+0.5878i$ \\
$4$  & Great Sea Water   & Renxu, Guihai & $4\pi/15$ & $0.6691+0.7431i$ \\
$5$  & Gold in the Sea   & Jiazi, Yichou & $\pi/3$ & $0.5+0.8660i$ \\
$6$  & Fire in the Furnace   & Bingyin, Dingmao & $2\pi/5$ & $0.3090+0.9511i$ \\
$7$  & Great Forest Wood   & Wuchen, Jisi & $7\pi/15$ & $0.1045+0.9945i$ \\
$8$  & Earth by the Roadside   & Gengwu, Xinwei & $8\pi/15$ & $-0.1045+0.9945i$ \\
$9$  & Metal of the Sword's Edge   & Renshen, Guiyou & $3\pi/5$ & $-0.3090+0.9511i$ \\
$10$ & Fire on the Mountain Top   & Jiaxu, Yihai & $2\pi/3$ & $-0.5+0.8660i$ \\
$11$ & Water in the Ravine   & Bingzi, Dingchou & $11\pi/15$ & $-0.6691+0.7431i$ \\
$12$ & Earth of the City Wall   & Wuyin, Jimao & $4\pi/5$ & $-0.8090+0.5878i$ \\
$13$ & White Wax Metal   & Gengchen, Xinsi & $13\pi/15$ & $-0.9135+0.4067i$ \\
$14$ & Willow Wood   & Renwu, Guiwei & $14\pi/15$ & $-0.9781+0.2079i$ \\
$15$ & Water in the Spring   & Jiashen, Yiyou & $\pi$ & $-1+0i$ \\
$16$ & Earth on the Roof   & Bingxu, Dinghai & $16\pi/15$ & $-0.9781-0.2079i$ \\
$17$ & Thunderbolt Fire   & Wuzi, Jichou & $17\pi/15$ & $-0.9135-0.4067i$ \\
$18$ & Pine and Cypress Wood   & Gengyin, Xinmao & $6\pi/5$ & $-0.8090-0.5878i$ \\
$19$ & Long Flowing Water   & Renchen, Guisi & $19\pi/15$ & $-0.6691-0.7431i$ \\
$20$ & Gold in the Sand   & Jiawu, Yiwei & $4\pi/3$ & $-0.5-0.8660i$ \\
$21$ & Fire at the Foot of the Mountain   & Bingshen, Dingyou & $7\pi/5$ & $-0.3090-0.9511i$ \\
$22$ & Wood of the Flat Ground   & Wuxu, Jihai & $22\pi/15$ & $-0.1045-0.9945i$ \\
$23$ & Earth on the Wall   & Gengzi, Xinchou & $23\pi/15$ & $0.1045-0.9945i$ \\
$24$ & Gold Foil   & Renyin, Guimao & $8\pi/5$ & $0.3090-0.9511i$ \\
$25$ & Fire Under the Lamp Cover   & Jiachen, Yisi & $5\pi/3$ & $0.5-0.8660i$ \\
$26$ & Water of the Heavenly River   & Bingwu, Dingwei & $26\pi/15$ & $0.6691-0.7431i$ \\
$27$ & Earth of the Great Post Station   & Wushen, Jiyou & $9\pi/5$ & $0.8090-0.5878i$ \\
$28$ & Gold of the Hairpin   & Gengxu, Xinhai & $28\pi/15$ & $0.9135-0.4067i$ \\
$29$ & Mulberry Wood   & Renzi, Guichou & $29\pi/15$ & $0.9781-0.2079i$ \\
\bottomrule
\end{tabular}
\end{center}

\subsection{\texorpdfstring{The sexagenary cycle ($T=60$)}{The sexagenary cycle (T=60)}}

\begin{definition}[Sexagenary-Cycle phase operators]\label{def:ganzhi_operators}
The stem--branch period is $T=60=4 \times 3 \times 5$, the bit width is $W=60$ and the bandwidth is $\Delta\theta=\pi/30$. The group structure is $C_{60} \cong C_4 \times C_3 \times C_5$ and the tensor shape is $4 \times 3 \times 5$. It is generated by the periodic convolution of the Heavenly Stems $\circledast$ the Earthly Branches.
\end{definition}

\begingroup
\centering
\begin{longtable}{>{\centering\arraybackslash}p{1cm}>{\centering\arraybackslash}p{1.5cm}>{\centering\arraybackslash}p{1.5cm}>{\centering\arraybackslash}p{2cm}>{\centering\arraybackslash}p{3.5cm}}
\toprule
$k$ & Stems and Branches & Five Phases & Phase $\theta_k$ & Operator $\mathcal{O}_k^{(60)}$ \\
\midrule
\endfirsthead

\multicolumn{5}{c}{\small(continued)} \\
\toprule
$k$ & Stems and Branches & Five Phases & Phase $\theta_k$ & Operator $\mathcal{O}_k^{(60)}$ \\
\midrule
\endhead

\midrule
\multicolumn{5}{r}{\small(continued on the next page)} \\
\endfoot

\bottomrule
\endlastfoot

$0$ & Jiayin & yang wood (beginning) & $0$ & $1+0i$ \\
$1$ & Yimao & yin wood & $\pi/30$ & $0.9945+0.1045i$ \\
$2$ & Bingchen & yang earth & $\pi/15$ & $0.9781+0.2079i$ \\
$3$ & Dingsi & yin fire & $\pi/10$ & $0.9511+0.3090i$ \\
$4$ & Wuwu & yang fire & $2\pi/15$ & $0.9135+0.4067i$ \\
$5$ & Jiwei & yin earth & $\pi/6$ & $0.8660+0.5i$ \\
$6$ & Gengshen & yang metal & $\pi/5$ & $0.8090+0.5878i$ \\
$7$ & Xinyou & yin metal & $7\pi/30$ & $0.7431+0.6691i$ \\
$8$ & Renxu & yang water & $4\pi/15$ & $0.6691+0.7431i$ \\
$9$ & Guihai & yin water & $3\pi/10$ & $0.5878+0.8090i$ \\
$10$ & Jiazi & yang wood & $\pi/3$ & $0.5+0.8660i$ \\
$11$ & Yichou & yin wood & $11\pi/30$ & $0.4067+0.9135i$ \\
$12$ & Bingyin & yang fire & $2\pi/5$ & $0.3090+0.9511i$ \\
$13$ & Dingmao & yin fire & $13\pi/30$ & $0.2079+0.9781i$ \\
$14$ & Wuchen & yang earth & $7\pi/15$ & $0.1045+0.9945i$ \\
$15$ & Jisi & yin earth & $\pi/2$ & $0+1i$ \\
$16$ & Gengwu & yang metal & $8\pi/15$ & $-0.1045+0.9945i$ \\
$17$ & Xinwei & yin metal & $17\pi/30$ & $-0.2079+0.9781i$ \\
$18$ & Renshen & yang water & $3\pi/5$ & $-0.3090+0.9511i$ \\
$19$ & Guiyou & yin water & $19\pi/30$ & $-0.4067+0.9135i$ \\
$20$ & Jiaxu & yang wood & $2\pi/3$ & $-0.5+0.8660i$ \\
$21$ & Yihai & yin wood & $7\pi/10$ & $-0.5878+0.8090i$ \\
$22$ & Bingzi & yang fire & $11\pi/15$ & $-0.6691+0.7431i$ \\
$23$ & Dingchou & yin fire & $23\pi/30$ & $-0.7431+0.6691i$ \\
$24$ & Wuyin & yang earth & $4\pi/5$ & $-0.8090+0.5878i$ \\
$25$ & Jimao & yin earth & $5\pi/6$ & $-0.8660+0.5i$ \\
$26$ & Gengchen & yang metal & $13\pi/15$ & $-0.9135+0.4067i$ \\
$27$ & Xinsi & yin metal & $9\pi/10$ & $-0.9511+0.3090i$ \\
$28$ & Renwu & yang water & $14\pi/15$ & $-0.9781+0.2079i$ \\
$29$ & Guiwei & yin water & $29\pi/30$ & $-0.9945+0.1045i$ \\
$30$ & Jiashen & yang wood & $\pi$ & $-1+0i$ \\
$31$ & Yiyou & yin wood & $31\pi/30$ & $-0.9945-0.1045i$ \\
$32$ & Bingxu & yang fire & $16\pi/15$ & $-0.9781-0.2079i$ \\
$33$ & Dinghai & yin fire & $11\pi/10$ & $-0.9511-0.3090i$ \\
$34$ & Wuzi & yang earth & $17\pi/15$ & $-0.9135-0.4067i$ \\
$35$ & Jichou & yin earth & $7\pi/6$ & $-0.8660-0.5i$ \\
$36$ & Gengyin & yang metal & $6\pi/5$ & $-0.8090-0.5878i$ \\
$37$ & Xinmao & yin metal & $37\pi/30$ & $-0.7431-0.6691i$ \\
$38$ & Renchen & yang water & $19\pi/15$ & $-0.6691-0.7431i$ \\
$39$ & Guisi & yin water & $13\pi/10$ & $-0.5878-0.9511i$ \\
$40$ & Jiawu & yang wood & $4\pi/3$ & $-0.5-0.8660i$ \\
$41$ & Yiwei & yin wood & $41\pi/30$ & $-0.4067-0.9135i$ \\
$42$ & Bingshen & yang fire & $7\pi/5$ & $-0.3090-0.9511i$ \\
$43$ & Dingyou & yin fire & $43\pi/30$ & $-0.2079-0.9781i$ \\
$44$ & Wuxu & yang earth & $22\pi/15$ & $-0.1045-0.9945i$ \\
$45$ & Jihai & yin earth & $3\pi/2$ & $0-1i$ \\
$46$ & Gengzi & yang metal & $23\pi/15$ & $0.1045-0.9945i$ \\
$47$ & Xinchou & yin metal & $47\pi/30$ & $0.2079-0.9781i$ \\
$48$ & Renyin & yang water & $8\pi/5$ & $0.3090-0.9511i$ \\
$49$ & Guimao & yin water & $49\pi/30$ & $0.4067-0.9135i$ \\
$50$ & Jiachen & yang wood & $5\pi/3$ & $0.5-0.8660i$ \\
$51$ & Yisi & yin wood & $17\pi/10$ & $0.5878-0.8090i$ \\
$52$ & Bingwu & yang fire & $26\pi/15$ & $0.6691-0.7431i$ \\
$53$ & Dingwei & yin fire & $53\pi/30$ & $0.7431-0.6691i$ \\
$54$ & Wushen & yang earth & $9\pi/5$ & $0.8090-0.5878i$ \\
$55$ & Jiyou & yin earth & $11\pi/6$ & $0.8660-0.5i$ \\
$56$ & Gengxu & yang metal & $28\pi/15$ & $0.9135-0.4067i$ \\
$57$ & Xinhai & yin metal & $19\pi/10$ & $0.9511-0.3090i$ \\
$58$ & Renzi & yang water & $29\pi/15$ & $0.9781-0.2079i$ \\
$59$ & Guichou & yin water (end) & $59\pi/30$ & $0.9945-0.1045i$ \\
\end{longtable}
\endgroup

\subsection{\texorpdfstring{The Sixty-Four Hexagrams ($T=64$)}{The Sixty-Four Hexagrams (T=64)}}

\begin{definition}[Sixty-Four-Hexagrams phase operators]\label{def:liushisigua_operators}
The period of the Sixty-Four Hexagrams is $T=64=8 \times 8$, the bit width is $W=64$ and the bandwidth is $\Delta\theta=\pi/32$. The group structure is $C_8 \times C_8$ and the tensor shape is $8 \times 8$. They are generated by nested epicycles of the Eight Trigrams $\times$ the Eight Trigrams.
\end{definition}

\begingroup
\small
\centering
\begin{longtable}{>{\centering\arraybackslash}p{0.8cm}>{\centering\arraybackslash}p{1.2cm}>{\centering\arraybackslash}p{1.2cm}>{\centering\arraybackslash}p{1.5cm}>{\centering\arraybackslash}p{1.5cm}>{\centering\arraybackslash}p{2.2cm}>{\centering\arraybackslash}p{2.8cm}}
\toprule
$k$ & Lower trigram & Upper trigram & Hexagram name & Five Phases & Phase $\theta_k$ & Operator $\mathcal{O}_k^{(64)}$ \\
\midrule
\endfirsthead

\multicolumn{7}{c}{\small(continued)} \\
\toprule
$k$ & Lower trigram & Upper trigram & Hexagram name & Five Phases & Phase $\theta_k$ & Operator $\mathcal{O}_k^{(64)}$ \\
\midrule
\endhead

\midrule
\multicolumn{7}{r}{\small(continued on the next page)} \\
\endfoot

\bottomrule
\endlastfoot

$0$ & Zhen ☳ & Zhen ☳ & Zhen (Thunder) & wood (beginning) & $0$ & $1+0i$ \\
$1$ & Xun ☴ & Zhen ☳ & Heng (Duration) & wood & $\pi/32$ & $0.9952+0.0980i$ \\
$2$ & Gen ☶ & Zhen ☳ & Xiaoguo (Small Preponderance) & earth & $\pi/16$ & $0.9808+0.1951i$ \\
$3$ & Li ☲ & Zhen ☳ & Feng (Abundance) & fire & $3\pi/32$ & $0.9569+0.2903i$ \\
$4$ & Qian ☰ & Zhen ☳ & Dazhuang (Great Power) & metal & $\pi/8$ & $0.9239+0.3827i$ \\
$5$ & Kun ☷ & Zhen ☳ & Yu (Enthusiasm) & earth & $5\pi/32$ & $0.8819+0.4714i$ \\
$6$ & Dui ☱ & Zhen ☳ & Guimei (The Marrying Maiden) & metal & $3\pi/16$ & $0.8315+0.5556i$ \\
$7$ & Kan ☵ & Zhen ☳ & Xie (Deliverance) & water & $7\pi/32$ & $0.7730+0.6344i$ \\
$8$ & Zhen ☳ & Xun ☴ & Yi (Increase) & wood & $\pi/4$ & $0.7071+0.7071i$ \\
$9$ & Xun ☴ & Xun ☴ & Xun (Wind) & wood & $9\pi/32$ & $0.6344+0.7730i$ \\
$10$ & Gen ☶ & Xun ☴ & Jian (Gradual Progress) & earth & $5\pi/16$ & $0.5556+0.8315i$ \\
$11$ & Li ☲ & Xun ☴ & Jiaren (The Family) & fire & $11\pi/32$ & $0.4714+0.8819i$ \\
$12$ & Qian ☰ & Xun ☴ & Xiaochu (The Taming Power of the Small) & metal & $3\pi/8$ & $0.3827+0.9239i$ \\
$13$ & Kun ☷ & Xun ☴ & Guan (Contemplation) & earth & $13\pi/32$ & $0.2903+0.9569i$ \\
$14$ & Dui ☱ & Xun ☴ & Zhongfu (Inner Truth) & metal & $7\pi/16$ & $0.1951+0.9808i$ \\
$15$ & Kan ☵ & Xun ☴ & Huan (Dispersion) & water & $15\pi/32$ & $0.0980+0.9952i$ \\
$16$ & Zhen ☳ & Gen ☶ & Yi (Nourishment) & wood & $\pi/2$ & $0+1i$ \\
$17$ & Xun ☴ & Gen ☶ & Gu (Work on What Has Been Spoiled) & wood & $17\pi/32$ & $-0.0980+0.9952i$ \\
$18$ & Gen ☶ & Gen ☶ & Gen (Mountain) & earth & $9\pi/16$ & $-0.1951+0.9808i$ \\
$19$ & Li ☲ & Gen ☶ & Bi (Grace) & fire & $19\pi/32$ & $-0.2903+0.9569i$ \\
$20$ & Qian ☰ & Gen ☶ & Dachu (The Taming Power of the Great) & metal & $5\pi/8$ & $-0.3827+0.9239i$ \\
$21$ & Kun ☷ & Gen ☶ & Bo (Splitting Apart) & earth & $21\pi/32$ & $-0.4714+0.8819i$ \\
$22$ & Dui ☱ & Gen ☶ & Sun (Decrease) & metal & $11\pi/16$ & $-0.5556+0.8315i$ \\
$23$ & Kan ☵ & Gen ☶ & Meng (Youthful Folly) & water & $23\pi/32$ & $-0.6344+0.7730i$ \\
$24$ & Zhen ☳ & Li ☲ & Shihe (Biting Through) & wood & $3\pi/4$ & $-0.7071+0.7071i$ \\
$25$ & Xun ☴ & Li ☲ & Ding (The Cauldron) & wood & $25\pi/32$ & $-0.7730+0.6344i$ \\
$26$ & Gen ☶ & Li ☲ & Lü (The Wanderer) & earth & $13\pi/16$ & $-0.8315+0.5556i$ \\
$27$ & Li ☲ & Li ☲ & Li (Fire) & fire & $27\pi/32$ & $-0.8819+0.4714i$ \\
$28$ & Qian ☰ & Li ☲ & Dayou (Great Possession) & metal & $7\pi/8$ & $-0.9239+0.3827i$ \\
$29$ & Kun ☷ & Li ☲ & Jin (Progress) & earth & $29\pi/32$ & $-0.9569+0.2903i$ \\
$30$ & Dui ☱ & Li ☲ & Kui (Opposition) & metal & $15\pi/16$ & $-0.9808+0.1951i$ \\
$31$ & Kan ☵ & Li ☲ & Weiji (Before Completion) & water & $31\pi/32$ & $-0.9952+0.0980i$ \\
$32$ & Zhen ☳ & Qian ☰ & Wuwang (Innocence) & wood & $\pi$ & $-1+0i$ \\
$33$ & Xun ☴ & Qian ☰ & Gou (Coming to Meet) & wood & $33\pi/32$ & $-0.9952-0.0980i$ \\
$34$ & Gen ☶ & Qian ☰ & Dun (Retreat) & earth & $17\pi/16$ & $-0.9808-0.1951i$ \\
$35$ & Li ☲ & Qian ☰ & Tongren (Fellowship with Men) & fire & $35\pi/32$ & $-0.9569-0.2903i$ \\
$36$ & Qian ☰ & Qian ☰ & Qian (Heaven) & metal & $9\pi/8$ & $-0.9239-0.3827i$ \\
$37$ & Kun ☷ & Qian ☰ & Pi (Standstill) & earth & $37\pi/32$ & $-0.8819-0.4714i$ \\
$38$ & Dui ☱ & Qian ☰ & Lü (Treading) & metal & $19\pi/16$ & $-0.8315-0.5556i$ \\
$39$ & Kan ☵ & Qian ☰ & Song (Conflict) & water & $39\pi/32$ & $-0.7730-0.6344i$ \\
$40$ & Zhen ☳ & Kun ☷ & Fu (Return) & wood & $5\pi/4$ & $-0.7071-0.7071i$ \\
$41$ & Xun ☴ & Kun ☷ & Sheng (Pushing Upward) & wood & $41\pi/32$ & $-0.6344-0.7730i$ \\
$42$ & Gen ☶ & Kun ☷ & Qian (Modesty) & earth & $21\pi/16$ & $-0.5556-0.8315i$ \\
$43$ & Li ☲ & Kun ☷ & Mingyi (Darkening of the Light) & fire & $43\pi/32$ & $-0.4714-0.8819i$ \\
$44$ & Qian ☰ & Kun ☷ & Tai (Peace) & metal & $11\pi/8$ & $-0.3827-0.9239i$ \\
$45$ & Kun ☷ & Kun ☷ & Kun (Earth) & earth & $45\pi/32$ & $-0.2903-0.9569i$ \\
$46$ & Dui ☱ & Kun ☷ & Lin (Approach) & metal & $23\pi/16$ & $-0.1951-0.9808i$ \\
$47$ & Kan ☵ & Kun ☷ & Shi (The Army) & water & $47\pi/32$ & $-0.0980-0.9952i$ \\
$48$ & Zhen ☳ & Dui ☱ & Sui (Following) & wood & $3\pi/2$ & $0-1i$ \\
$49$ & Xun ☴ & Dui ☱ & Daguo (Great Preponderance) & wood & $49\pi/32$ & $0.0980-0.9952i$ \\
$50$ & Gen ☶ & Dui ☱ & Xian (Influence) & earth & $25\pi/16$ & $0.1951-0.9808i$ \\
$51$ & Li ☲ & Dui ☱ & Ge (Revolution) & fire & $51\pi/32$ & $0.2903-0.9569i$ \\
$52$ & Qian ☰ & Dui ☱ & Guai (Breakthrough) & metal & $13\pi/8$ & $0.3827-0.9239i$ \\
$53$ & Kun ☷ & Dui ☱ & Cui (Gathering Together) & earth & $53\pi/32$ & $0.4714-0.8819i$ \\
$54$ & Dui ☱ & Dui ☱ & Dui (Lake) & metal & $27\pi/16$ & $0.5556-0.8315i$ \\
$55$ & Kan ☵ & Dui ☱ & Kun (Oppression) & water & $55\pi/32$ & $0.6344-0.7730i$ \\
$56$ & Zhen ☳ & Kan ☵ & Zhun (Difficulty at the Beginning) & wood & $7\pi/4$ & $0.7071-0.7071i$ \\
$57$ & Xun ☴ & Kan ☵ & Jing (The Well) & wood & $57\pi/32$ & $0.7730-0.6344i$ \\
$58$ & Gen ☶ & Kan ☵ & Jian (Obstruction) & earth & $29\pi/16$ & $0.8315-0.5556i$ \\
$59$ & Li ☲ & Kan ☵ & Jiji (After Completion) & fire & $59\pi/32$ & $0.8819-0.4714i$ \\
$60$ & Qian ☰ & Kan ☵ & Xu (Waiting) & metal & $15\pi/8$ & $0.9239-0.3827i$ \\
$61$ & Kun ☷ & Kan ☵ & Bi (Holding Together) & earth & $61\pi/32$ & $0.9569-0.2903i$ \\
$62$ & Dui ☱ & Kan ☵ & Jie (Limitation) & metal & $31\pi/16$ & $0.9808-0.1951i$ \\
$63$ & Kan ☵ & Kan ☵ & Kan (Water) & water (end) & $63\pi/32$ & $0.9952-0.0980i$ \\
\end{longtable}
\endgroup

\section{Application of periodic convolution in the I Ching symbol system}

\section{Master table of the tensor structure of the periodic systems}

\begin{center}
\begin{tabular}{>{\centering\arraybackslash}p{2.2cm}>{\centering\arraybackslash}p{1cm}>{\centering\arraybackslash}p{2.5cm}>{\centering\arraybackslash}p{3cm}>{\centering\arraybackslash}p{2cm}>{\centering\arraybackslash}p{2.5cm}}
\toprule
Symbol system & $T$ & Prime factorisation & Group structure & Tensor shape & Nested generation \\
\midrule
Taiji & $1$ & $1$ & $C_1$ & $1$ & basic \\
Yin-Yang & $2$ & $2$ & $C_2$ & $2$ & basic \\
Sancai & $3$ & $3$ & $C_3$ & $3$ & basic \\
the Four Symbols & $4$ & $2^2$ & $C_4$ & $4$ & Yin-Yang $\times$ Yin-Yang \\
the Five Phases & $5$ & $5$ & $C_5$ & $5$ & basic \\
the Seven Luminaries & $7$ & $7$ & $C_7$ & $7$ & basic \\
the Eight Trigrams & $8$ & $2^3$ & $C_8$ & $8$ & Yin-Yang $\times$ the Four Symbols \\
the Heavenly Stems & $10$ & $2 \times 5$ & $C_2 \times C_5$ & $2 \times 5$ & Yin-Yang $\times$ the Five Phases \\
the Earthly Branches & $12$ & $4 \times 3$ & $C_4 \times C_3$ & $4 \times 3$ & the Four Symbols $\times$ Sancai \\
the twenty-four solar terms & $24$ & $8 \times 3$ & $C_8 \times C_3$ & $8 \times 3$ & the Eight Trigrams $\times$ Sancai \\
the thirty Na-Yin & $30$ & $2 \times 3 \times 5$ & $C_2 \times C_3 \times C_5$ & $2 \times 3 \times 5$ & Yin-Yang $\times$ Sancai $\times$ the Five Phases \\
the sexagenary cycle & $60$ & $4 \times 3 \times 5$ & $C_4 \times C_3 \times C_5$ & $4 \times 3 \times 5$ & the Heavenly Stems $\circledast$ the Earthly Branches \\
the Sixty-Four Hexagrams & $64$ & $8 \times 8$ & $C_8 \times C_8$ & $8 \times 8$ & the Eight Trigrams $\times$ the Eight Trigrams \\
\bottomrule
\end{tabular}
\end{center}

\begin{remark}[Mathematical positioning of the two tables]
The `table of phase operators of the thirteen symbol systems' is the period-by-period instantiation of definition \ref{def:PPO_yijing} and definition \ref{def:tensor_representation_yijing}: the correspondence of the phase operator $\mathcal{O}_k^{(T)} = e^{i2\pi k/T}$ of each $T$ with its symbol, its Five-Phase assignment and its tensor index is the complete element table of the cyclic group $C_T$ of that period. The `master table of the tensor structure of the periodic systems' collects the prime factorisation, the group structure, the tensor shape and the source of nested generation of the thirteen symbol systems, and is a systematic index of theorem \ref{thm:cyclic_direct_product_yijing} and proposition \ref{prop:irrep_tensor_product_yijing}. Neither table carries an independent theorem; their mathematical positioning is to render the group-theoretic and representation-theoretic conclusions of the main text as a symbol correspondence that can be looked up in a table.
\end{remark}

\section{The major-luck cycle: approximate replication of the stems and branches on the ten-year scale}
\label{sec:dayun_yijing}

\begin{remark}[Positioning of this section: connection to the five-step interface of the `Periodic Systems' chapter]\label{rem:dayun_positioning_yijing}
The major-luck cycle (\textit{dayun}) is a ten-year-per-step sequence used in traditional fate calculation to characterise an individual's long-term fortunes. This section does not treat it as a new astronomical period; rather it characterises it as an \textbf{approximate replication, on the ten-year scale, of the sexagenary structure already established in this chapter}. Specifically it connects to the five-step interface of the `Periodic Systems' chapter:
\begin{enumerate}[label=(\arabic*)]
    \item \textbf{Decision}: reduce the major-luck sequence to a discrete dynamical system on $C_{60}$, cite \cref{thm:ps_periodicity_decision} to decide its period type;
    \item \textbf{Construction}: from the major-luck stem--branch sequence construct the amplitude series, the frequency series and the phase series, and assemble them into the major-luck epicycle group;
    \item \textbf{Adjustment}: dynamically adjust the major-luck epicycle group by the starting age, the choice of direction and the decrease of the residual energy;
    \item \textbf{Convolution}: introduce amplitude convolution, phase convolution and frequency convolution on the major-luck epicycle group, and connect them with the periodic convolution framework of this chapter;
    \item \textbf{Integration}: connect to the tool chains of the three chapters `Discrete Sampling and Intercalation of Periodic Signals', `Harmonic Analysis' and `Time Series Analysis', forming a major-luck periodicity diagnostic procedure.
\end{enumerate}
This section uses tools already established in this chapter: the periodic phase operator $\mathcal O_k^{(T)}$ (\cref{def:PPO_yijing}), the nested epicycle generation rule (\cref{def:epicycle_nesting_rule}) and the periodic-convolution generation of the sexagenary cycle (\cref{thm:ganzhi_periodic_convolution}). It also cites the conclusions of the calendar chapter on the solar-term phase (\cref{def:jieqi_zhongqi_dingyue}) and on fixed length (\cref{def:fixed_length}).
\end{remark}

\begin{remark}[Guide to the seven themes of this section]\label{rem:dayun_seven_themes}
Translation---the ordering advance of the major-luck sequence along the sexagenary cycle (\cref{def:dayun_sequence_yijing});
Scaling---the tenfold replication of scale in the major-luck cycle ($10$ years per step, \cref{thm:dayun_three_periods_yijing});
Rotation---the direction $\sigma=\pm 1$ and the precession of the relative phase with respect to the current year (\cref{prop:dayun_liunian_phase_yijing});
Conservation---the threefold conservation of the year-stem yin-yang $\epsilon$, the step length $L$ and the month pillar $M_0$ (\cref{thm:dayun_conservation_yijing});
Symmetry---the directional duality between yang males and yin females and between yin males and yang females (\cref{prop:dayun_direction_duality_yijing});
Closure---the closure of the major-luck sequence with the period $60$ steps ($600$ years) (\cref{thm:dayun_periodicity_yijing});
Period---the ten-year step, the hundred-year Heavenly-Stem period, the hundred-and-twenty-year Earthly-Branch period and the six-hundred-year stem--branch period (\cref{thm:dayun_three_periods_yijing}).
\end{remark}

\subsection{The stems and branches as a basic mechanism of rational approximation}

\begin{definition}[Heavenly-Stem period and Earthly-Branch period]\label{def:stem_branch_period_yijing}
The \textbf{Heavenly Stems} are ten symbols, with period $T_{\text{stem}}=10$;
the \textbf{Earthly Branches} are twelve symbols, with period $T_{\text{branch}}=12$.
The \textbf{stems and branches} are the synchronous pairing of the Heavenly Stems with the Earthly Branches, and their period is
\[
T_{\text{stem-branch}}=\operatorname{lcm}(T_{\text{stem}},T_{\text{branch}})=\operatorname{lcm}(10,12)=60.
\]
It is written as the cyclic group $C_{60}=C_{10}\times C_{12}/\sim$, where the equivalence relation $\sim$ is given by the parity-synchronisation condition $(a,b)\sim(a',b')$ if and only if $a-b\equiv a'-b'\pmod 2$.
\end{definition}

\begin{proposition}[The stems and branches as the fibre product of $C_{10}$ and $C_{12}$]\label{prop:ganzhi_fibre_product_yijing}
The cyclic group $C_{60}$ is isomorphic to the fibre product of $C_{10}$ and $C_{12}$ under the parity projections:
\[
C_{60}\cong\{(a,b)\in C_{10}\times C_{12}\mid a\equiv b\pmod 2\}.
\]
\end{proposition}

\begin{proof}
By the generalisation of the Chinese remainder theorem, $C_{10}\times C_{12}$ has order $120$. The parity projections $C_{10}\to C_2$ and $C_{12}\to C_2$ are both surjective; the fibre product therefore has size $|C_{10}\times C_{12}|/|C_2|=60$. The map from $C_{60}$ to the fibre product, $n\mapsto(n\bmod 10,n\bmod 12)$, is a group homomorphism and a bijection, hence an isomorphism.
\end{proof}

\begin{definition}[Empirical values of the astronomical periods]\label{def:planetary_periods_yijing}
The following values of the astronomical periods come from modern observation (in years):
\[
\begin{aligned}
T_J&\approx 11.862 &&(\text{Jupiter's orbital period}),\\
T_S&\approx 29.457 &&(\text{Saturn's orbital period}),\\
T_{JS}&=\frac{T_J T_S}{T_S-T_J}\approx 19.859 &&(\text{Jupiter--Saturn conjunction}).
\end{aligned}
\]
\end{definition}

\begin{proposition}[Approximation of the year stems and branches to Jupiter]\label{prop:year_ganzhi_jupiter_yijing}
The $60$-year stem--branch period corresponds to about $5$ revolutions of Jupiter:
\[
\frac{60}{T_J}=\frac{60}{11.862}\approx 5.058\approx 5,
\]
with a relative error of about $1.2\%$. This approximation is the structural origin of the \textbf{Jupiter year-reckoning system} (in which Jupiter takes about $12$ years to complete one circuit of the heavens): the Earthly-Branch period $12$ is an integer approximation to the period of Jupiter.
\end{proposition}

\begin{proof}
$60/11.862=5.058$, whose relative error with respect to the integer $5$ is $(5.058-5)/5=0.0116\approx 1.2\%$. The relative error between the Earthly-Branch period $12$ and the period of Jupiter $T_J\approx 11.862$ is $1.1\%$, of the same order.
\end{proof}

\begin{proposition}[Approximation of the year stems and branches to the Jupiter--Saturn conjunction]\label{prop:year_ganzhi_jupiter_saturn_yijing}
The $60$-year stem--branch period corresponds to about $3$ Jupiter--Saturn conjunctions:
\[
\frac{60}{T_{JS}}=\frac{60}{19.859}\approx 3.021\approx 3,
\]
with a relative error of about $0.7\%$.
\end{proposition}

\begin{proof}
$60/19.859=3.021$, whose relative error with respect to the integer $3$ is $0.021/3\approx 0.7\%$.
\end{proof}

\begin{remark}[The year stems and branches as a double approximation]\label{rem:ganzhi_double_approximation_yijing}
\cref{prop:year_ganzhi_jupiter_yijing} and \cref{prop:year_ganzhi_jupiter_saturn_yijing} reveal that the $60$-year stem--branch period is a double rational approximation: to \textbf{$5$ revolutions of Jupiter} and to \textbf{$3$ Jupiter--Saturn conjunctions}. These two approximations correspond exactly to the composition of $T_{\text{stem}}=10$ and $T_{\text{branch}}=12$:
\begin{itemize}
    \item $T_{\text{branch}}=12\approx T_J$ (the period of Jupiter), and Jupiter completes about $5$ revolutions within $60$ years;
    \item $T_{\text{stem-branch}}=60\approx 3\,T_{JS}$ ($3$ Jupiter--Saturn conjunctions).
\end{itemize}
The role of the Heavenly-Stem period $T_{\text{stem}}=10$ is to modulate the single period $T_{\text{branch}}=12$, so that the fibre product of $C_{10}\times C_{12}$ yields a composite rhythm of $60$ years. This structure corresponds algebraically, one to one, with \cref{thm:ganzhi_periodic_convolution} of this chapter (the periodic-convolution generation of the sexagenary cycle).
\end{remark}

\subsection{The constituents of the major-luck cycle and its tenfold replication of scale}

\begin{definition}[Major-luck target and decision input]\label{def:dayun_target_yijing}
A \textbf{major-luck target} is characterised by a quadruple:
\[
\mathcal D_{\mathrm{D}}=(M_0,\sigma,a_0,L),\qquad L=10\ \text{years},
\]
where $M_0\in C_{60}$ is the month pillar (the starting point), $\sigma\in\{+1,-1\}$ is the direction, $a_0\ge 0$ is the starting age of the cycle and $L=10$ is the step length. The state space is the cyclic group $\mathcal S=C_{60}$, and the evolution map is
\[
\Phi:\mathcal S\to\mathcal S,\qquad \Phi(m)=m+\sigma\pmod{60}.
\]
\end{definition}

\begin{definition}[Major-luck sequence]\label{def:dayun_sequence_yijing}
Let the decision input of the major-luck target be $\mathcal D_{\mathrm{D}}=(M_0,\sigma,a_0,L)$. The \textbf{major-luck sequence} is defined by
\[
M_k=M_0+k\sigma\pmod{60},\qquad k=0,1,2,\dots
\]
The $k$-th major-luck step corresponds to the time interval $[a_0+kL,\ a_0+(k+1)L)$.
\end{definition}

\begin{theorem}[The three fundamental periods of the major-luck cycle]\label{thm:dayun_three_periods_yijing}
The major-luck sequence has three fundamental periods:
\begin{enumerate}[label=(\roman*)]
    \item \textbf{Heavenly-Stem period}: $T_{\text{stem}}^{\mathrm{(D)}}=T_{\text{stem}}\cdot L=10\cdot 10=100$ years;
    \item \textbf{Earthly-Branch period}: $T_{\text{branch}}^{\mathrm{(D)}}=T_{\text{branch}}\cdot L=12\cdot 10=120$ years;
    \item \textbf{Stem--branch period}: $T_{\text{stem-branch}}^{\mathrm{(D)}}=T_{\text{stem-branch}}\cdot L=60\cdot 10=600$ years.
\end{enumerate}
Here the step length $L=10$ years is exactly equal to the Heavenly-Stem period, which makes the major-luck cycle a \textbf{replication of the year stems and branches on the ten-year scale}.
\end{theorem}

\begin{proof}
The major-luck Heavenly Stems have $10$ positions, each corresponding to one step, so the number of steps needed for the Heavenly Stems to return to their place is $10$, that is, $10\times 10=100$ years. Likewise, the Earthly Branches need $12$ steps to return to their place, that is, $120$ years. The stems and branches together need $\operatorname{lcm}(10,12)=60$ steps to return to their place, that is, $600$ years.
\end{proof}

\begin{proposition}[Periodicity decision theorem for the major-luck cycle]\label{thm:dayun_periodicity_decision_yijing}
Let the decision map of the major-luck target be $\Phi(m)=m+\sigma\pmod{60}$. Then:
\begin{enumerate}[label=(\Alph*)]
    \item \textbf{Sufficient decision}: $\Phi^{60}=\mathrm{id}$, so the target is a uniformly periodic major-luck system;
    \item \textbf{Equivalence of the decision}: $\operatorname{ord}(\Phi)=60$;
    \item \textbf{Pointwise criterion}: the orbit of the major-luck sequence at $M_0$ has length $60$;
    \item \textbf{Phase independence}: the starting age $a_0$ changes only the starting position of the sequence, not the period.
\end{enumerate}
\end{proposition}

\begin{proof}
(A) $\Phi^{60}(m)=m+60\sigma\equiv m\pmod{60}$.
(B) If $\operatorname{ord}(\Phi)=k<60$, then $60\mid k$ (since $\gcd(\sigma,60)=1$), a contradiction.
(C) By (B) and \cref{thm:finite_order_orbit} of the `Operations and Periodicity' chapter.
(D) By \cref{def:dayun_sequence_yijing}, $M_k$ does not depend on $a_0$.
\end{proof}

\begin{corollary}[The triple astronomical approximation of the major-luck cycle]\label{cor:dayun_triple_approximation_yijing}
The three fundamental periods of the major-luck cycle approximate the following astronomical rhythms respectively:
\begin{enumerate}[label=(\roman*)]
    \item \textbf{the Heavenly-Stem period of $100$ years} $\approx$ $5$ Jupiter--Saturn conjunctions:
    \[
    \frac{100}{T_{JS}}=\frac{100}{19.859}\approx 5.036\approx 5,\quad\text{relative error }0.7\%;
    \]
    \item \textbf{the Earthly-Branch period of $120$ years} $\approx$ $10$ revolutions of Jupiter:
    \[
    \frac{120}{T_J}=\frac{120}{11.862}\approx 10.12\approx 10,\quad\text{relative error }1.2\%;
    \]
    \item \textbf{the stem--branch period of $600$ years} $\approx$ $50$ revolutions of Jupiter $\approx$ $30$ Jupiter--Saturn conjunctions:
    \[
    \frac{600}{T_J}\approx 50.6\approx 50,\qquad \frac{600}{T_{JS}}\approx 30.2\approx 30,
    \]
    with relative errors of $1.2\%$ and $0.7\%$ respectively.
\end{enumerate}
\end{corollary}

\begin{proof}
Term by term:
(i) $100/19.859=5.036$, error $0.036/5=0.7\%$.
(ii) $120/11.862=10.12$, error $0.12/10=1.2\%$.
(iii) $600/11.862=50.58$, error $0.58/50=1.2\%$; $600/19.859=30.21$, error $0.21/30=0.7\%$.
\end{proof}

\begin{theorem}[Duality of scale between the major-luck cycle and the year stems and branches]\label{thm:dayun_vs_year_ganzhi_yijing}
The major-luck cycle and the year stems and branches form a duality of scale:
\begin{center}
\begin{tabular}{lccc}
\toprule
Scale & Heavenly-Stem period & Earthly-Branch period & Stem--branch period \\
\midrule
Year stems and branches & $10$ & $12$ & $60$ \\
Major-luck stems and branches & $100$ & $120$ & $600$ \\
\midrule
Astronomical approximation of the year stems and branches & Jupiter--Saturn conjunctions $\times 3$ ($60$ years) & Jupiter $\times 5$ ($60$ years) & as at left \\
Astronomical approximation of the major-luck cycle & Jupiter--Saturn conjunctions $\times 5$ ($100$ years) & Jupiter $\times 10$ ($120$ years) & Jupiter $\times 50$ ($600$ years) \\
\bottomrule
\end{tabular}
\end{center}
The two columns are completely parallel in structure: each $10$-year step of the major-luck cycle is precisely a replication of the Heavenly-Stem period $T_{\text{stem}}=10$ of the year stems and branches. This structure has the same origin as the periodic-convolution generation of the sexagenary cycle in \cref{thm:ganzhi_periodic_convolution} of this chapter.
\end{theorem}

\begin{proof}
The periods of the year stems and branches are $(T_{\text{stem}},T_{\text{branch}},T_{\text{stem-branch}})=(10,12,60)$, see \cref{def:stem_branch_period_yijing}. The major-luck cycle unfolds with step length $L=T_{\text{stem}}=10$ years, so each of the three periods is multiplied by $10$, giving $(100,120,600)$. The year stems and branches approximate $(T_J\times 5,\ T_{JS}\times 3)=(59.3,59.6)$, whilst the major-luck stems and branches approximate $(T_J\times 10,\ T_{JS}\times 5)=(118.6,99.3)$; the two are of the same order of accuracy (a relative error of the order of $1\%$).
\end{proof}

\begin{example}[Numerical verification of the $120$-year major-luck Earthly-Branch period against $10$ revolutions of Jupiter]\label{ex:dayun_jupiter_120_yijing}
$10\times T_J = 10\times 11.862 = 118.62$ years. The difference between the $120$-year major-luck Earthly-Branch period and $118.62$ years is $1.38$ years, a relative error of
\[
\frac{1.38}{120}\approx 1.15\%.
\]
This is of exactly the same order as the relative error of $1.2\%$ with which the $12$-year year-stem-and-branch Earthly-Branch period approximates Jupiter's $11.862$ years.
\end{example}

\begin{example}[Numerical verification of the $100$-year major-luck Heavenly-Stem period against $5$ Jupiter--Saturn conjunctions]\label{ex:dayun_js_100_yijing}
$5\times T_{JS}=5\times 19.859=99.30$ years. The difference between the $100$-year major-luck Heavenly-Stem period and $99.30$ years is $0.70$ years, a relative error of
\[
\frac{0.70}{100}=0.70\%.
\]
\end{example}

\subsection{The binary character of the major-luck direction}

\begin{definition}[Year-stem yin-yang and sex]\label{def:year_stem_yin_yang_yijing}
Let the index of the year-pillar Heavenly Stem be $\epsilon_{\text{stem}}\in\{0,1,\dots,9\}$ (Jia $=0$, Yi $=1$, $\dots$, Gui $=9$). Define the \textbf{yin-yang marker} of the year stem as
\[
\epsilon=\begin{cases}1,&\epsilon_{\text{stem}}\in\{0,2,4,6,8\}\ (\text{yang}),\\0,&\epsilon_{\text{stem}}\in\{1,3,5,7,9\}\ (\text{yin}).\end{cases}
\]
Define the \textbf{sex marker} $\eta\in\{0,1\}$ (male $=1$, female $=0$).
\end{definition}

\begin{definition}[Major-luck direction]\label{def:dayun_direction_yijing}
The \textbf{direction} $\sigma\in\{+1,-1\}$ of the major-luck cycle is determined by $(\epsilon,\eta)$:
\[
\sigma=2\cdot\mathbb 1_{\{\epsilon=\eta\}}-1=
\begin{cases}+1,&\epsilon=\eta\ (\text{yang male, yin female; forward ordering}),\\-1,&\epsilon\neq\eta\ (\text{yin male, yang female; reverse ordering}).\end{cases}
\]
\end{definition}

\begin{proposition}[Duality of the major-luck direction]\label{prop:dayun_direction_duality_yijing}
The major-luck direction satisfies:
\begin{enumerate}[label=(\roman*)]
    \item $\sigma(\epsilon,1-\eta)=-\sigma(\epsilon,\eta)$;
    \item $\sigma(1-\epsilon,\eta)=-\sigma(\epsilon,\eta)$;
    \item $\sigma(1-\epsilon,1-\eta)=\sigma(\epsilon,\eta)$.
\end{enumerate}
\end{proposition}

\begin{proof}
(i) If $\epsilon=\eta$, then $\epsilon\neq 1-\eta$, giving $\sigma=-1=-\sigma(\epsilon,\eta)$; if $\epsilon\neq\eta$, then $\epsilon=1-\eta$, giving $\sigma=+1=-\sigma(\epsilon,\eta)$. (ii) and (iii) are analogous.
\end{proof}

\begin{corollary}[Equivalence classes of the major-luck direction]\label{cor:dayun_direction_classes_yijing}
The four cases $(\epsilon,\eta)\in\{0,1\}^2$ fall under $\sigma$ into two equivalence classes:
\[
\{(1,1),(0,0)\}\xrightarrow{\sigma}+1,\qquad \{(1,0),(0,1)\}\xrightarrow{\sigma}-1.
\]
The two equivalence classes are each closed under the composite duality $(\epsilon,\eta)\mapsto(1-\epsilon,1-\eta)$.
\end{corollary}

\begin{proof}
This follows directly from \cref{def:dayun_direction_yijing} and \cref{prop:dayun_direction_duality_yijing}(iii).
\end{proof}

\subsection{Quantification of the starting age of the cycle}

\begin{definition}[Starting age of the cycle]\label{def:qiyun_age_yijing}
Let an individual be born at the instant $t_{\text{birth}}$. Let $\Delta d$ be
\[
\Delta d=\begin{cases}t_{\text{next}}-t_{\text{birth}},&\sigma=+1,\\t_{\text{birth}}-t_{\text{prev}},&\sigma=-1,\end{cases}
\]
where $t_{\text{prev}}$ and $t_{\text{next}}$ are the instants of the nearest solar terms before and after the birth (for the definition of a solar term see \cref{def:jieqi_zhongqi_dingyue} of the calendar chapter). The \textbf{starting age of the cycle} is
\[
a_0=\frac{\Delta d}{3}\quad(\text{years}),
\]
according to the traditional quantification rule $3$ days $=1$ year, $1$ day $=4$ months, $1$ double-hour $=10$ days.
\end{definition}

\begin{proposition}[Quantitative bound on the starting age of the cycle]\label{prop:qiyun_age_bound_yijing}
Let the time interval between two adjacent solar terms be $\Delta T_{\text{jie}}$. Then
\[
0\le a_0\le \frac{\Delta T_{\text{jie}}}{3}.
\]
Under the mean-solar-term method $\Delta T_{\text{jie}}=T_{\text{trop}}/24\approx 15.2184$ days (\cref{def:fixed_length} of the calendar chapter), hence $a_0\le 5.07$ years.
\end{proposition}

\begin{proof}
By \cref{def:qiyun_age_yijing}, $0\le\Delta d\le\Delta T_{\text{jie}}$. Substituting into $a_0=\Delta d/3$ gives the result.
\end{proof}

\begin{remark}[Phase interpretation of the starting age of the cycle]\label{rem:qiyun_phase_yijing}
The starting age $a_0=\Delta d/3$ is a \textbf{phase offset}---it measures the relative position of the instant of birth within the solar-term cycle. The value $a_0$ does not change the periodicity of the major-luck sequence (\cref{thm:dayun_periodicity_decision_yijing}(D)); it changes only the translation of the starting point of the sequence along the time axis. This is consistent with the phase-encoding framework of the periodic phase operator $\mathcal O_k^{(T)}=e^{i2\pi k/T}$ of this chapter (\cref{def:PPO_yijing}).
\end{remark}

\begin{example}[Numerical computation of the starting age of the cycle]\label{ex:qiyun_calc_yijing}
Let an individual be born $9$ days after Lichun, with $\sigma=+1$. Then $a_0=9/3=3$ years, and the major-luck sequence begins at the age of $3$.
\end{example}

\subsection{The periodic structure of the major-luck sequence and its astronomical meaning}

\begin{theorem}[Periodic structure of the major-luck sequence]\label{thm:dayun_periodicity_yijing}
Let the major-luck sequence be $M_k=M_0+k\sigma\pmod{60}$. Then:
\begin{enumerate}[label=(\roman*)]
    \item \textbf{Algebraic period}: $M_{k+60}=M_k$ holds for all $k\in\mathbb Z$;
    \item \textbf{Time period}: $60\times 10=600$ years;
    \item \textbf{Single orbit}: $C_{60}$ consists of a single orbit under the major-luck translation group $\langle\Phi\rangle$;
    \item \textbf{No conserved quantities}: every invariant on $C_{60}$ with respect to the major-luck translation must be constant.
\end{enumerate}
\end{theorem}

\begin{proof}
(i) By $\Phi^{60}=\mathrm{id}$ (\cref{thm:dayun_periodicity_decision_yijing}(A)).
(ii) Each step is $10$ years, so $60$ steps are $600$ years.
(iii) $C_{60}$ consists of a single orbit under $\Phi$.
(iv) By \cref{thm:period_not_imply_conservation} of the `Operations and Periodicity' chapter, in a single-orbit system a conserved quantity must be constant.
\end{proof}

\begin{remark}[The $600$-year period as the joint locking of the astronomical approximations]\label{rem:dayun_600_year_yijing}
By \cref{thm:dayun_three_periods_yijing} and \cref{cor:dayun_triple_approximation_yijing}, the $600$-year major-luck period simultaneously approximates:
\begin{itemize}
    \item $50$ revolutions of Jupiter: $50\times 11.862=593.1$ years, a relative error of $1.2\%$;
    \item $20$ revolutions of Saturn: $20\times 29.457=589.14$ years, a relative error of $1.8\%$;
    \item $30$ Jupiter--Saturn conjunctions: $30\times 19.859=595.77$ years, a relative error of $0.7\%$.
\end{itemize}
These three astronomical rhythms are approximated simultaneously on the $600$-year scale---this is the deeper astronomical meaning of the periodic structure of the major-luck cycle. It explains why the seemingly artificial construction of $60$ steps $\times$ $10$ years possesses long-term stability and empirical value within the traditional calendrical system.
\end{remark}

\begin{example}[An instance of the periodic closure of the major-luck sequence]\label{ex:dayun_period_closure_yijing}
Let $M_0=$ `Jiayin' (index $0$) and $\sigma=+1$. The $60$-th major-luck step is $M_{60}=M_0+60\equiv 0\pmod{60}$, returning to Jiayin. The sequence closes after $600$ years. This is consistent with the periodic closure of the sexagenary cycle of this chapter (\cref{thm:ganzhi_periodic_convolution}).
\end{example}

\subsection{The relative phase between the major-luck cycle and the current year}

\begin{definition}[Current year]\label{def:liunian_yijing}
The \textbf{current year} (\textit{liunian}) refers to the year pillar of a specific year. Let the starting age of the cycle be $a_0$; then the current year in year $t$ ($t\ge a_0$) is
\[
N_t=N_{a_0}+(t-a_0)\pmod{60}.
\]
\end{definition}

\begin{proposition}[Relative phase between the major-luck cycle and the current year]\label{prop:dayun_liunian_phase_yijing}
Let $N_{a_0}$ be the year pillar of the year in which the cycle starts and $M_0$ the month pillar. Then the relative phase between the $k$-th major-luck step and the current year is
\[
\Delta_k=N_{a_0+kL}-M_k\equiv (N_{a_0}-M_0)+k(L-\sigma)\pmod{60},
\]
where $L=10$. In particular:
\begin{enumerate}[label=(\roman*)]
    \item \textbf{Forward ordering} ($\sigma=+1$): $\Delta_k\equiv\Delta_0+9k\pmod{60}$;
    \item \textbf{Reverse ordering} ($\sigma=-1$): $\Delta_k\equiv\Delta_0+11k\pmod{60}$.
\end{enumerate}
\end{proposition}

\begin{proof}
The starting point of the $k$-th major-luck step corresponds to the age $a_0+kL$, and the current year is $N_{a_0+kL}=N_{a_0}+kL\pmod{60}$. Subtraction gives $\Delta_k\equiv N_{a_0}-M_0+k(L-\sigma)\pmod{60}$. Substituting $L=10$ and $\sigma=\pm 1$ gives the result.
\end{proof}

\begin{corollary}[Period of the relative phase between the major-luck cycle and the current year]\label{cor:dayun_liunian_period_yijing}
The period of the relative phase $\Delta_k$ is $T_\Delta=60/\gcd(L-\sigma,60)$. For $\sigma=+1$ ($L-\sigma=9$, $\gcd(9,60)=3$), $T_\Delta=20$ steps; for $\sigma=-1$ ($L-\sigma=11$, $\gcd(11,60)=1$), $T_\Delta=60$ steps.
\end{corollary}

\begin{proof}
By \cref{prop:dayun_liunian_phase_yijing}, $\Delta_k$ is the translation sequence on $C_{60}$ with step length $(L-\sigma)$, whose period is $60/\gcd(L-\sigma,60)$ (\cref{lem:order_cyclic} of the `Operations and Periodicity' chapter).
\end{proof}

\begin{remark}[Asymmetry of the phase periods of forward and reverse ordering]\label{rem:dayun_asymmetry_yijing}
The major-luck--current-year phase period is $20$ steps ($200$ years) for forward ordering and $60$ steps ($600$ years) for reverse ordering. This asymmetry arises from $\gcd(9,60)=3$ whilst $\gcd(11,60)=1$---$9$ and $60$ have a common factor, whereas $11$ and $60$ are coprime. It explains the structural difference between forward and reverse ordering in traditional current-year coupling.
\end{remark}

\begin{example}[A concrete computation of the closure of the major-luck--current-year phase]\label{ex:dayun_liunian_closure_yijing}
Let $M_0=$ `Jiayin' (index $0$), $N_{a_0}=$ `Bingyin' (index $12$), $\sigma=+1$ and $a_0=3$. Then
\[
\Delta_0\equiv 2-50\equiv 12\pmod{60}.
\]
The relative phase increases in steps of $9$: $\Delta_1=21$, $\Delta_2=30$, $\Delta_3=39$, $\Delta_4=48$, $\Delta_5=57$, $\dots$. By \cref{cor:dayun_liunian_period_yijing}, $T_\Delta=20$, that is, the phase closes after $20$ steps ($200$ years).
\end{example}

\subsection{Conserved quantities of the major-luck system}

\begin{theorem}[The three invariants of the major-luck system]\label{thm:dayun_conservation_yijing}
The major-luck system has three invariants:
\begin{enumerate}[label=(\roman*)]
    \item \textbf{Conservation of direction}: $\sigma\in\{+1,-1\}$ is unchanged throughout the major-luck sequence;
    \item \textbf{Conservation of step length}: $L=10$ years is unchanged throughout the major-luck sequence;
    \item \textbf{Conservation of the starting point}: $M_0\in C_{60}$ is the anchor of the sequence and can be uniquely recovered from it.
\end{enumerate}
\end{theorem}

\begin{proof}
(i) By \cref{def:dayun_direction_yijing}, $\sigma$ is uniquely determined by $(\epsilon,\eta)$.
(ii) By \cref{def:dayun_sequence_yijing}, $L=10$ is a fixed parameter.
(iii) Given $(\sigma,a_0)$ and the sequence $\{M_k\}$, the value $M_0$ is read off directly at $k=0$.
\end{proof}

\subsection{Construction and adjustment of the major-luck epicycle group}

\begin{definition}[Major-luck amplitude series, frequency series and phase series]\label{def:dayun_series_yijing}
Let $\{m_k\}$ be the indices of the major-luck sequence on $C_{60}$. If
\[
m_k=\sum_{j}A_j e^{i(\omega_j k+\phi_j)}+\varepsilon_k,
\]
then $\mathcal A_m(k)=\sum_j A_j e^{ijk}$, $\mathcal F_m(s)=\sum_j\omega_j s^j$ and $\mathcal P_m(k)=\sum_j e^{i(jk+\phi_j)}$ are called the amplitude series, the frequency series and the phase series of the major-luck sequence.
\end{definition}

\begin{definition}[Major-luck epicycle group]\label{def:dayun_epicycle_group_yijing}
The \textbf{epicycle group} of the major-luck target is
\[
\mathfrak E_{\mathrm{D}}=\{(A_j,\omega_j,\phi_j)\}_{j\in\mathbb Z},
\]
and its superposition is the major-luck period series $m_k=\sum_j A_j e^{i(\omega_j k+\phi_j)}$. This formulation is consistent with the nested epicycle generation rule of this chapter (\cref{def:epicycle_nesting_rule}).
\end{definition}

\begin{proposition}[Fourier expansion of the major-luck sequence]\label{prop:dayun_fourier_yijing}
Let the major-luck sequence take index values on $C_{60}$. Then its $60$-point DFT is
\[
\hat m_j=\frac{1}{60}\sum_{k=0}^{59}m_k e^{-i2\pi jk/60},\qquad j=0,1,\dots,59,
\]
and $m_k=\sum_{j=0}^{59}\hat m_j e^{i2\pi jk/60}$.
\end{proposition}

\begin{proof}
By \cref{thm:dayun_periodicity_yijing}(i), $\{m_k\}$ is a finite sequence of period $60$. By the standard result on discrete Fourier series, the $60$-point DFT gives the above decomposition.
\end{proof}

\begin{definition}[Dynamic adjustment of the major-luck epicycle group]\label{def:dayun_adaptive_yijing}
Let the current epicycle group be $\mathfrak E^{(j)}$ and let the residual be $r^{(j)}=m-m^{(j)}$. The dynamic adjustment consists of:
\begin{enumerate}[label=(\roman*)]
    \item \textbf{Cardinality adjustment}: if $\|r^{(j)}\|^2>\delta$, add one epicycle according to the peak of the residual periodogram;
    \item \textbf{Parameter adjustment}: update the existing parameters by regularised least squares.
\end{enumerate}
\end{definition}

\begin{proposition}[Monotonicity and convergence of the major-luck dynamic adjustment]\label{prop:dayun_adaptive_convergence_yijing}
After adding the dominant frequency component of the residual as in \cref{def:dayun_adaptive_yijing}(i), the residual energy is monotonically non-increasing; if the target is a \textbf{real-valued finite periodic sequence} of length $T_0$ ($T_0\in\mathbb N$) with $\delta=0$ (here $T_0=60$), then after at most $\lfloor 60/2\rfloor+1=31$ steps the dynamic adjustment leaves a zero residual (for a \textbf{complex-valued} sequence the corresponding number of steps is at most $60$).
\end{proposition}

\begin{proof}
By the orthogonal decomposition, $\|r^{(j)}\|^2=\|r^{(j+1)}\|^2+\|\text{the removed component}\|^2$. The DFT of a \textbf{real-valued} finite periodic sequence of length $60$ is complete on $31$ harmonics (including the constant term); a \textbf{complex-valued} sequence requires $60$ components.
\end{proof}

\subsection{Major-luck convolution}

\begin{definition}[Major-luck amplitude convolution, phase convolution and frequency convolution]\label{def:dayun_convolution_yijing}
Let two major-luck epicycle groups be $\mathfrak E=\{(A_n,\omega_n,\phi_n)\}$ and $\mathfrak E'=\{(B_m,\nu_m,\psi_m)\}$. Define:
\[
\mathbb A_{mn}=A_nB_m,\quad \mathbb P_{mn}=e^{i(\phi_n-\psi_m)},\quad \mathbb F_{mn}=\omega_n\nu_m,\quad \mathbb T_{mn}=A_nB_m e^{i(\phi_n-\psi_m)}.
\]
\end{definition}

\begin{proposition}[Phase convolution of the major-luck cycle and the current year]\label{prop:dayun_liunian_convolution_yijing}
The phase convolution $\mathbb P_{kt}=e^{i(\phi_k^M-\phi_t^N)}$ of the major-luck sequence with the current-year sequence gives the relative phase of the two; its argument is precisely the phase corresponding to $\Delta_k$ and measures the strength of the major-luck--current-year coupling. This is consistent with the periodic convolution framework of this chapter (\cref{def:PPO_periodic_convolution}).
\end{proposition}

\begin{proof}
This follows directly from \cref{def:dayun_convolution_yijing} and \cref{prop:dayun_liunian_phase_yijing}.
\end{proof}

\subsection{Integration with discrete sampling, harmonic analysis and time series analysis}

\begin{definition}[Major-luck periodicity diagnostic procedure]\label{def:dayun_diagnosis_yijing}
A given major-luck system is diagnosed as follows:
\begin{enumerate}[label=(\roman*)]
    \item \textbf{Periodicity decision}: verify that the major-luck sequence has period $60$ steps (\cref{thm:dayun_periodicity_yijing});
    \item \textbf{Direction diagnosis}: compute $\sigma$ from $(\epsilon,\eta)$ (\cref{def:dayun_direction_yijing});
    \item \textbf{Starting-age diagnosis}: compute $a_0$ from $\Delta d/3$ (\cref{def:qiyun_age_yijing});
    \item \textbf{Phase diagnosis}: compute the major-luck--current-year relative phase $\Delta_k$ (\cref{prop:dayun_liunian_phase_yijing});
    \item \textbf{Astronomical-approximation diagnosis}: verify the approximation of Jupiter and of the Jupiter--Saturn conjunction by the three periods $100$, $120$ and $600$ years (\cref{cor:dayun_triple_approximation_yijing}).
\end{enumerate}
\end{definition}

\begin{theorem}[The major-luck periodicity diagnostic procedure and the interface with three chapters]\label{thm:dayun_diagnosis_integration_yijing}
The five steps of \cref{def:dayun_diagnosis_yijing} connect respectively to:
\begin{enumerate}[label=(\roman*)]
    \item the chapter `Discrete Sampling and Intercalation of Periodic Signals', \cref{def:uniform_sampling_parallel}, \cref{thm:nyquist_shannon_parallel} (the quantification of the starting age is a sampling from continuous time to discrete years and months);
    \item the chapter `Harmonic Analysis', \cref{def:fourier_series}, \cref{thm:parseval_fourier} (the DFT decomposition of the major-luck sequence);
    \item the chapter `Time Series Analysis', \cref{def:ts_period_detection}, \cref{prop:ts_period_significance} (the significance test of the major-luck--current-year phase peak).
\end{enumerate}
\end{theorem}

\begin{proof}
(i) The starting age $a_0=\Delta d/3$ is a sampling from continuous time to discrete years and months. (ii) The major-luck sequence has period $60$, and its DFT belongs to the tools of harmonic analysis. (iii) The major-luck--current-year relative phase forms a spectral peak within the period, whose significance can be decided by the periodogram and Fisher's test.
\end{proof}

\begin{remark}[The periodic-mathematical essence of the major-luck system]\label{rem:dayun_math_essence_yijing}
The periodic-mathematical essence of the major-luck system is summarised as follows:
\begin{enumerate}[label=(\roman*)]
    \item \textbf{Starting point} $M_0$: the month pillar, determined jointly by the solar-term phase and the year stem (\cref{def:jieqi_zhongqi_dingyue}, \cref{def:year_stem_yin_yang_yijing} of the calendar chapter);
    \item \textbf{Direction} $\sigma$: determined by the binary duality of the year-stem yin-yang $\epsilon$ and the sex $\eta$ (\cref{def:dayun_direction_yijing});
    \item \textbf{Phase offset} $a_0$: the starting age of the cycle, given by the number of days from the birth to the solar term according to the quantification rule $3$ days $=1$ year (\cref{def:qiyun_age_yijing});
    \item \textbf{Sequence} $M_k=M_0+k\sigma$: a translation on the sexagenary cycle, of period $60$ steps.
\end{enumerate}
The major-luck cycle is not a new astronomical period, but the \textbf{replication, on the ten-year scale, of the rational approximation of the year stems and branches to Jupiter and to the Jupiter--Saturn conjunction}: the year stems and branches approximate $5$ revolutions of Jupiter within $60$ years, and the major-luck stems and branches approximate $50$ revolutions of Jupiter within $600$ years, the two relative errors being of the same order (of the order of $1\%$).
\end{remark}

\begin{remark}[Connections between the major-luck cycle and the other periodic systems of this chapter]\label{rem:dayun_relation_to_yijing}
The relation of the major-luck cycle to the other periodic systems of this chapter may be summarised as four connections:
\begin{enumerate}[label=(\roman*)]
    \item \textbf{With the sexagenary cycle}: the major-luck sequence is a translation orbit of the sexagenary cycle (\cref{thm:ganzhi_periodic_convolution}) on $C_{60}$---it introduces no new algebraic object, but unfolds the cyclic structure of the sexagenary cycle afresh on a new time scale;
    \item \textbf{With the periodic phase operators}: the $k$-th major-luck step corresponds to the phase encoding of the periodic phase operator $\mathcal O_k^{(60)}$ (\cref{def:PPO_yijing}), and the major-luck direction $\sigma$ corresponds to the sign of the phase rotation;
    \item \textbf{With periodic convolution}: the relative phase of the major-luck cycle and the current year is the phase convolution of the major-luck epicycle group with the current-year epicycle group (\cref{prop:dayun_liunian_convolution_yijing});
    \item \textbf{With the calendar system}: the quantification rule for the starting age ($3$ days $=1$ year) and intercalation in the calendar chapter (\cref{thm:calendar_continued_fraction_best}) both belong to the framework of best rational approximation, but their objects differ---the former quantifies the mapping from time to years and months, whilst the latter quantifies the integer pairs of the lunisolar calendar.
\end{enumerate}

\textbf{Core insight}: the major-luck cycle is the replication, on the ten-year scale, of the rational approximation of the year stems and branches to Jupiter and to the Jupiter--Saturn conjunction. The year stems and branches approximate $(T_J\times 5,\ T_{JS}\times 3)$ by $(10,12,60)$, whilst the major-luck stems and branches approximate $(T_J\times 50,\ T_{JS}\times 30)$ by $(100,120,600)$, and the accuracy of approximation is of the same order at the two scales (of the order of $1\%$). The core contribution of the major-luck cycle is this: it implements the sexagenary algebraic structure established in this chapter, turning it from a `symbol system' into a `verifiable sequence of astronomical approximations'---the step length $L=T_{\text{stem}}=10$ years is precisely a replication of the Heavenly-Stem period of the year stems and branches.
\end{remark}

\section*{Seven-theme comparison table for the major-luck cycle}

\begin{center}
\renewcommand{\arraystretch}{1.3}
\begin{tabular}{>{\centering\arraybackslash}p{2.2cm}>{\centering\arraybackslash}p{10.5cm}}
\toprule
Theme & Corresponding content in the major-luck cycle \\
\midrule
Translation & the ordering advance of the major-luck sequence $M_k=M_0+k\sigma$ along the sexagenary cycle (\cref{def:dayun_sequence_yijing}) \\
Scaling & the tenfold replication of scale in the major-luck cycle: $T_{\text{stem}}^{\mathrm{(D)}}=100$ years, $T_{\text{branch}}^{\mathrm{(D)}}=120$ years, $T_{\text{stem-branch}}^{\mathrm{(D)}}=600$ years (\cref{thm:dayun_three_periods_yijing}) \\
Rotation & the major-luck direction $\sigma=\pm 1$ and the precession of the relative phase $\Delta_k$ with respect to the current year (\cref{prop:dayun_liunian_phase_yijing}) \\
Conservation & the threefold conservation of the direction $\sigma$, the step length $L$ and the starting point $M_0$ (\cref{thm:dayun_conservation_yijing}) \\
Symmetry & the directional duality between yang males and yin females and between yin males and yang females (\cref{prop:dayun_direction_duality_yijing}, \cref{cor:dayun_direction_classes_yijing}) \\
Closure & the closure of the major-luck sequence with the period $60$ steps ($600$ years) (\cref{thm:dayun_periodicity_yijing}) \\
Period & the triple astronomical approximation: $100$ years $\approx$ $5$ Jupiter--Saturn conjunctions, $120$ years $\approx$ $10$ revolutions of Jupiter, $600$ years $\approx$ $50$ revolutions of Jupiter (\cref{cor:dayun_triple_approximation_yijing}) \\
\bottomrule
\end{tabular}
\end{center}
\begin{remark}[Unity with the whole book]\label{rem:yijing_book_interface}
The periodic structure of the I Ching system established in this chapter is a concrete realisation of periodic mathematics in a traditional symbol system. The core observation is that \textbf{a phase is an operator}: every phase $\mathcal{O}_k^{(T)} = e^{i2\pi k/T}$ of every period is a relation operator. The yang operator ($\mathcal{O}_0 = 1$) defines the identity relation, the yin operator ($\mathcal{O}_{T/2} = -1$) defines the opposition relation, and the remaining operators define partial rotation relations.

The nested epicycle generation rule---for every single step forward of the outer circle, the inner circle rotates through one full turn---unifies the process by which a basic period generates a large period: Yin-Yang $\times$ the Five Phases $=$ the Heavenly Stems ($T=10$), the Four Symbols $\times$ Sancai $=$ the Earthly Branches ($T=12$), and the Eight Trigrams $\times$ the Eight Trigrams $=$ the Sixty-Four Hexagrams ($T=64$). The sexagenary cycle is generated from the synchronous advance of the Heavenly Stems and the Earthly Branches by periodic convolution (\cref{thm:ganzhi_periodic_convolution}). The He Tu complex function establishes the Fourier grading $\mathbb{Z}_{10} \cong \mathbb{Z}_2 \times \mathbb{Z}_5$ (\cref{thm:hetu_fourier_decomposition_yijing}), and the Lo Shu matrix establishes the translation--rotation spectral decomposition of the $3 \times 3$ magic square (\cref{thm:luoshu_decomposition_yijing}). The two connect the traditional symbol system with standard mathematical tools from the two directions of complex analysis and linear algebra respectively.

\textbf{The place of the major-luck cycle in the algebraic structure of this chapter}: the major-luck cycle is not a new algebraic object, but a \textbf{time-series realisation} of the sexagenary cycle $C_{60}$ on the ten-year scale. It implements the abstract phase operators established in this chapter (\cref{def:PPO_yijing}) and the periodic convolution framework (\cref{thm:ganzhi_periodic_convolution}) as a verifiable sequence of astronomical approximations: the year stems and branches approximate $5$ revolutions of Jupiter and $3$ Jupiter--Saturn conjunctions by $(10,12,60)$ (\cref{prop:year_ganzhi_jupiter_yijing}, \cref{prop:year_ganzhi_jupiter_saturn_yijing}), whilst the major-luck stems and branches approximate $50$ revolutions of Jupiter and $30$ Jupiter--Saturn conjunctions by $(100,120,600)$ (\cref{cor:dayun_triple_approximation_yijing}), and the accuracy of approximation is of the same order at the two scales (of the order of $1\%$). This structure shows that the sexagenary cycle is not only a symbol system but also a \textbf{sequence of rational approximations} that can be tested against astronomical observation at several scales---$10$, $60$, $120$ and $600$ years. The algebraic framework of this chapter (phase operators, epicycle nesting, periodic convolution) and the time-series realisation of the major-luck cycle together constitute the two sides---algebraic structure and astronomical content---of the sixty-year period.
\end{remark}

\section*{Seven-theme comparison table for this chapter}

\begin{center}
\renewcommand{\arraystretch}{1.3}
\begin{tabular}{>{\centering\arraybackslash}p{2.2cm}>{\centering\arraybackslash}p{10.5cm}}
\toprule
Theme & Corresponding content in this chapter \\
\midrule
Translation & the ordering advance of the Heavenly-Stem ($T=10$) and Earthly-Branch ($T=12$) sequences, generating the sexagenary cycle through parity synchronisation $n \equiv m \pmod 2$ (\cref{ch20_ganzhi_parity_filter}); the ordering advance of the major-luck sequence along the sexagenary cycle (\cref{def:dayun_sequence_yijing}) \\
Scaling & the doubling hierarchy of hexagram levels, the $2^n$ nesting: Yin-Yang $2^1 \to$ the Four Symbols $2^2 \to$ the Eight Trigrams $2^3$, and the Eight Trigrams $\times$ the Eight Trigrams $= 8 \times 8 = 64$ (\cref{ex:liushisigua_generation}); the tenfold replication of scale in the major-luck cycle (\cref{thm:dayun_three_periods_yijing}) \\
Rotation & the rotation of phase points by $2\pi k/T$ by the phase operator $\mathcal{O}_k^{(T)} = e^{i2\pi k/T}$ (\cref{def:PPO_yijing}, \cref{prop:PPO_relation_structure}); the Lo Shu rotational subspace (\cref{thm:luoshu_decomposition_yijing}); the major-luck direction $\sigma=\pm 1$ and the precession of the relative phase $\Delta_k$ with respect to the current year (\cref{prop:dayun_liunian_phase_yijing}) \\
Balance & the combinatorial conservation of the generation from the Eight Trigrams to the Sixty-Four Hexagrams: the number of independent phases equals the product of the bit widths $T_{\text{inner}} \times T_{\text{outer}}$ (\cref{thm:epicycle_nesting_generation}); the antipodal balance sum of the Lo Shu matrix is $10$ (\cref{prop:luoshu_antipodal_yijing}); the threefold conservation of the major-luck direction $\sigma$, the step length $L$ and the starting point $M_0$ (\cref{thm:dayun_conservation_yijing}) \\
Symmetry & the Yin-Yang involution symmetry: the yin operator $\mathcal{O}_{T/2}^{(T)} = -1$ gives the opposite vector when multiplied by $-1$, the $C_2$ involution (\cref{prop:PPO_relation_structure}); the antipodal antisymmetry of the Lo Shu matrix (\cref{prop:luoshu_antipodal_yijing}); the directional duality between yang males and yin females and between yin males and yang females (\cref{prop:dayun_direction_duality_yijing}) \\
Closure & the cyclic closure by which the sexagenary cycle ($T=60$) and the Sixty-Four Hexagrams ($T=64$) each return to their starting point (\cref{thm:ganzhi_periodic_convolution}, \cref{ex:liushisigua_generation}); the closure of the major-luck sequence with the period $60$ steps ($600$ years) (\cref{thm:dayun_periodicity_yijing}) \\
Period & the periodic systems of the thirteen symbol systems, such as the Heavenly Stems $T=10$, the Earthly Branches $T=12$ and the Eight Trigrams $T=8$; the Fourier grading of the He Tu complex function (\cref{thm:hetu_fourier_decomposition_yijing}); the rotation period $\pi/\sqrt{6}$ of the Lo Shu matrix (\cref{cor:luoshu_structure}); the triple astronomical approximation of the major-luck cycle (\cref{cor:dayun_triple_approximation_yijing}) \\
\bottomrule
\end{tabular}
\end{center}

\section*{Exercises}

\begin{enumerate}[leftmargin=3em, label=\textbf{Exercise \arabic*}]
    \item Prove the \textbf{periodicity decision theorem for I Ching targets} of this chapter (\ref{thm:yijing_periodicity_decision}). Hint: first write down the decision map and the evolution group of the target of this chapter, then apply the proof structure of the general criterion of the `Periodic Systems' chapter.

    \item Verify that the hypotheses of the \textbf{periodicity decision theorem for the major-luck cycle} (\ref{thm:dayun_periodicity_decision_yijing}) hold for the domain objects of this chapter, and give a counterexample in which the hypotheses fail.

    \item Examine the behaviour of the \textbf{monotonicity of the cardinality adjustment} (\ref{prop:yijing_adaptive_growth}) under degenerate parameters, and explain how the conclusion degenerates accordingly.

    \item Explain the relation between the \textbf{correspondence of the three parameters: radius, phase and angular speed} (\ref{prop:yijing_epicycle_correspondence}) and the corresponding general conclusion of the `Periodic Systems' chapter (whether it agrees with it, is a special case of it, or requires additional conditions).

    \item Using the seven-theme comparison table for this chapter, point out the concrete carrier of each of the three themes \textbf{translation}, \textbf{scaling} and \textbf{rotation} in this chapter, and give a domain-data example for each.

    \item Synthesis: restate the diagnostic procedure of this chapter using the five-step standard interface of the `Periodic Systems' chapter (table \ref{tab:ps_five_steps}), noting one by one the theorem of this chapter on which each step depends: \ref{thm:yijing_periodicity_decision}, \ref{thm:dayun_periodicity_decision_yijing}, \ref{prop:yijing_adaptive_growth}, \ref{prop:yijing_epicycle_correspondence}.
\end{enumerate}


